\appendix

\newcommand{\G}{{\mathscr G}}
\newcommand{\V}{{\mathscr V}}
\newcommand{\bG}{\mathbf{G}}
\newcommand{\bH}{\mathbf{H}}
\newcommand{\bR}{\mathbf{R}}
\newcommand{\bC}{\mathbf{C}}
\newcommand{\bP}{\mathbf{P}}
\newcommand{\Proj}{{\mathbb P}}
\renewcommand{\P}{{\mathbb P}}
\newcommand{\PPP}{\mathscr P}
\newcommand{\MMM}{\mathscr M}
\newcommand{\WWW}{\mathscr W}
\newcommand{\sA}{\mathscr A}
\newcommand{\sC}{\mathscr C}
\newcommand{\sO}{\mathscr O}
\newcommand{\sV}{\mathscr V}
\renewcommand{\sS}{\mathscr S}
\newcommand{\fg}{\mathfrak g}
\newcommand{\aaa}{\mathbf{a}}
\renewcommand{\H}{\mathscr H}
\newcommand{\fh}{\mathfrak h}
\renewcommand{\emptyset}{\varnothing}
\newcommand{\cusp}{\text{cusp}}
\newcommand{\bv}{\mathbf{v}}
\newcommand{\bu}{\mathbf{u}}
\newcommand{\borser}[1]{#1^{\text{BS}}}
\newcommand{\closure}[1]{(#1)^{-}}
\renewcommand{\norm}[1]{\|#1\|}
\newcommand{\Vor}{Vorono\v{\i}\xspace}
\newcommand{\vv}{{\mathbf {v}}}
\renewcommand{\aa}{{\mathbf {a}}}
\newcommand{\bb}{{\mathbf {b}}}
\newcommand{\ww}{{\mathbf {w}}}
\newcommand{\xx}{{\mathbf {x}}}
\newcommand{\uu}{{\mathbf {u}}}
\newcommand{\yy}{{\mathbf {y}}}
\newcommand{\qq}{{\mathbf {q}}}
\newcommand{\sets}[1]{[\![#1]\!]}
\newcommand{\coinv}[1]{(S_{#1})_{\Gamma }}
\renewcommand{\riso}{\xrightarrow{\sim}}


\chapter{Computing in Higher Rank}\label{ch:appendix}
\noindent{\Large\bf by Paul E. Gunnells}

%%%%%%%%%%%%%%%%%%
%                %
%  introduction  %
%                %
%%%%%%%%%%%%%%%%%%
\section{Introduction}\label{introduction}
%
%\subsection*{}
This book has addressed the theoretical and practical problems of
performing computations with modular forms.  Modular forms are the
simplest examples of the general theory of automorphic forms attached
to a reductive algebraic group $G$ with an arithmetic
subgroup~$\Gamma$; they are the case $G=\SL_{2} (\R)$ with $\Gamma$ a
congruence subgroup of $\SL_{2} (\Z)$.  For such pairs $(G,\Gamma)$
the Langlands philosophy asserts that there should be deep connections
between automorphic forms and arithmetic, connections that are
revealed through the action of the Hecke operators on spaces of
automorphic forms.  There have been many profound advances in recent
years in our understanding of these phenomena, for example:
\begin{itemize}
\item the establishment of the modularity of
elliptic curves defined over $\Q$ \cite{wiles:fermat, taylor-wiles:fermat, diamond,
cdt, breuil-conrad-diamond-taylor}, 
\item the proof by Harris--Taylor of the local Langlands
correspondence \cite{harris.taylor}, and  
\item Lafforgue's proof of the
global Langlands correspondence for function fields \cite{lafforgue}.
\end{itemize}
Nevertheless, we are still far from seeing that the links between
automorphic forms and arithmetic hold in the broad scope in which 
they are generally believed.  Hence one has the natural problem of
studying spaces of automorphic forms computationally.

%\subsection*{}
%
The goal of this appendix is to describe some computational techniques
for automorphic forms.  We focus on the case $G= \SL_{n} (\R)$ and
$\Gamma \subset \SL_{n} (\Z)$, since the automorphic forms that arise
are one natural generalization of modular forms, and since this is the
setting for which we have the most tools available.  In fact, we do not
work directly with automorphic forms, but rather with the cohomology
of the arithmetic group $\Gamma$ with certain coefficient modules.
This is the most natural generalization of the tools developed in
previous chapters.

Here is a brief overview of the contents.  Section \ref{aut.forms}
gives background on automorphic forms and the cohomology of arithmetic
groups and explains why the two are related.  In Section
\ref{comb.models} we describe the basic topological tools used to
compute the cohomology of $\Gamma$ explicitly.  Section
\ref{mod.symb.section} defines the Hecke operators, describes the
generalization of the modular symbols from Chapter \ref{ch:modsym} to
higher rank, and explains how to compute the action of the Hecke
operators on the top degree cohomology group.  Section
\ref{other.coho} discusses computation of the Hecke action on
cohomology groups below the top degree.  Finally, Section \ref{apps}
briefly discusses some related material and presents some open
problems.

\subsection{}
%
The theory of automorphic forms is notorious for the difficulty of its
prerequisites.  Even if one is only interested in the cohomology of
arithmetic groups---a small part of the full theory---one needs
considerable background in algebraic groups, algebraic topology, and
representation theory.  This is somewhat reflected in our
presentation, which falls far short of being self-contained.  Indeed,
a complete account would require a long book of its own.  We have
chosen to sketch the foundational material and to provide many
pointers to the literature; good general references are
\cite{borel.wallach, harder.notes, schwermer.article, vogan}.  We hope
that the energetic reader will follow the references and fill many
gaps on his/her own.

The choice of topics presented here is heavily influenced (as usual)
by the author's interests and expertise.  There are many computational
topics in the cohomology of arithmetic groups we have completely
omitted, including the trace formula in its many incarnations
\cite{gross.pollack}, the explicit Jacquet--Langlands correspondence
\cite{dembele,whitehouse}, and moduli space techniques
\cite{faber-van.der.geer, van.der.geer}.  We encourage the reader to investigate
these extremely interesting and useful techniques.

\subsection{Acknowledgements}
I thank Avner Ash, John Cremona, Mark McConnell, and Dan Yasaki for helpful
comments.  I also thank the NSF for support.  

\section{Automorphic Forms and Arithmetic Groups}\label{aut.forms}
%
%\subsection{Modular forms and cohomology}
\subsection{}
%
Let $\Gamma = \Gamma_{0} (N)\subset \SL_{2} (\Z)$ be the usual Hecke
congruence subgroup of matrices upper-triangular mod $N$. Let
$Y_{0} (N)$ be the modular curve $\Gamma \backslash \fh$, and let
$X_{0} (N)$ be its canonical compactification obtained by adjoining
cusps.  For any integer $k\geq 2$, let $S_{k} (N)$ be the space of
weight $k$ holomorphic cuspidal modular forms on $\Gamma$.  According
to Eichler--Shimura \cite[Chapter 8]{shimura:intro}, we have the isomorphism
\begin{equation}\label{eq:es}
H^{1} (X_{0} (N); \C)\riso S_{2} (N)\oplus \overline{S_{2} (N)},
\end{equation} 
where the bar denotes complex conjugation and where the isomorphism
is one of Hecke modules. 

More generally, for any integer $n\geq 0$, let $M_{n}\subset \C [x,y]$
be the subspace of degree $n$ homogeneous polynomials.  The space
$M_{n}$ admits a representation of $\Gamma $ by the ``change of
variables'' map
\begin{equation}\label{eq:rep}
\left(\begin{array}{cc}
a&b\\
c&d
\end{array} \right)\cdot p (x,y) = p (dx-by, -cx+ay).
\end{equation}
This induces a local system $\widetilde{M}_{n}$ on the curve $X_{0}
(N)$.\footnote{The classic references for cohomology with local
systems are \cite[Section~31]{steenrod} and \cite[Ch. V]{eilenberg}.  A more
recent exposition (in the language of \v Cech cohomology and locally
constant sheaves) can be found in \cite[II.13]{bott.tu}.  For an
exposition tailored to our needs, 
see \cite[Section~2.9]{harder.notes}.}  Then the
analogue of \eqref{eq:es} for higher-weight modular forms is the
isomorphism
\begin{equation}\label{eq:eshw}
H^{1} (X_{0} (N); \widetilde{M}_{k-2})\riso S_{k} (N)\oplus
\overline{S_{k} (N)}.  
\end{equation}
Note that \eqref{eq:eshw} reduces to \eqref{eq:es} when $k=2$.

Similar considerations apply if we work with the open curve $Y_{0}
(N)$ instead, except that Eisenstein series also contribute to the
cohomology.  More precisely, let $\Eis_{k} (N)$ be the space of weight
$k$ Eisenstein series on $\Gamma_{0} (N)$.  Then \eqref{eq:eshw}
becomes
\begin{equation}\label{eq:eseis}
H^{1} (Y_{0} (N); \widetilde{M}_{k-2})\riso S_{k} (N)\oplus
\overline{S_{k} (N)}\oplus \Eis_{k} (N).  
\end{equation}
These isomorphisms lie at the heart of the modular symbols method.

%\subsection{Modular forms and functions on $\SL_{2}
%(\R)$}\label{ss:pre.aut.forms}
\subsection{}\label{ss:pre.aut.forms}
%
The first step on the path to general automorphic forms is a
reinterpretation of modular forms in terms of functions on $\SL_{2}
(\R)$.  Let $\Gamma \subset \SL_{2}
(\Z)$ be a congruence subgroup.  A weight $k$ modular form on $\Gamma$
is a holomorphic function $f\colon \fh \rightarrow \C$ satisfying
the transformation property
\[
f ((az+b)/ (cz+d)) = j (\gamma, z)^{k}f (z), \quad \gamma = \left(\begin{array}{cc}
a&b\\
c&d
\end{array} \right)\in \Gamma, \quad z \in \fh.
\]  
Here $j (\gamma , z)$ is the \defn{automorphy factor} $cz+d$.  There
are some additional conditions $f$ must satisfy at the cusps of $\fh
$, but these are not so important for our discussion.

The group $G=\SL_{2} (\R)$ acts transitively on $\fh$, with the
subgroup $K=\SO (2)$ fixing $i$.  Thus $\fh$ can be written as the
quotient $G/K$.  From this, we see that $f$ can be viewed as a function
$G\rightarrow \C$ that is \emph{$K$-invariant on the right} and that
satisfies a certain symmetry condition with respect to the
\emph{$\Gamma$-action on the left}.  Of course not every $f$ with these
properties is a modular form: some extra data is needed to take the
role of holomorphicity and to handle the behavior at the cusps.
Again, this can be ignored right now.

We can turn this interpretation around as follows.  Suppose
$\varphi$ is a function $G\rightarrow \C$ that is \defn{$\Gamma
$-invariant on the left}, that is, $\varphi (\gamma g) = \varphi (g)$
for all $\gamma \in \Gamma$.  Hence $\varphi$ can be thought of as a
function $\varphi \colon \Gamma \backslash G\rightarrow \C$.  We
further suppose that $\varphi$ satisfies a certain symmetry condition
with respect to the \emph{$K$-action on the right.}  In particular,
any matrix $m\in K$ can be written 
\begin{equation}\label{eq:m-theta}
m = \left(\begin{array}{cc}
\cos \theta & -\sin \theta \\
\sin \theta & \cos \theta \\
\end{array} \right), \quad \theta \in \R ,
\end{equation}
with $\theta$ uniquely determined modulo $2\pi$.  Let $\zeta_{m}$ be
the complex number $e^{i\theta}$.  Then the $K$-symmetry we require is 
\[
\varphi (gm) = \zeta_{m}^{-k}\varphi(z), \quad m\in K,
\]
where $k$ is some fixed nonnegative integer.

It turns out that such functions $\varphi$ are very closely related to
modular forms: \emph{any} $f\in S_{k} (\Gamma)$ uniquely determines such a
function $\varphi_{f} \colon \Gamma \backslash G \rightarrow \C$.  The
correspondence is very simple.  Given a weight $k$ modular form $f$, define
\begin{equation}\label{eq:corresp}
\varphi_{f} (g) := f (g\cdot i)j (g,i)^{-k}.
\end{equation}
We claim $\varphi_{f}$ is left $\Gamma$-invariant and satisfies the 
desired $K$-symmetry on the right.
Indeed, since $j$ satisfies the
cocycle property $$j (gh, z) = j (g,h\cdot z)j (h,z),$$ we have 
\[
\varphi_{f} (\gamma g) = f
((\gamma g)\cdot i)j (\gamma g,i)^{-k} = j (\gamma ,g\cdot i)^{k}f (g\cdot
i)j (\gamma ,g\cdot i)^{-k}j (g,i)^{-k} = \varphi_{f} (g).
\]
Moreover, any
$m\in K$ stabilizes $i$.  Hence 
\[
\varphi_{f} (gm) = f ((gm)\cdot i) j
(gm,i)^{-k} = f (g\cdot i) j (m,i)^{-k}j (g,m\cdot i)^{-k}.
\]
From
\eqref{eq:m-theta} we have $j (m,i)^{-k} = (\cos \theta + i\sin
\theta)^{-k} = \zeta_{m}^{-k}$, and thus $\varphi_{f} (gm) =
 \zeta^{-k}_{m}\varphi_{f} (g)$.

Hence in \eqref{eq:corresp} the weight and the automorphy factor
``untwist'' the $\Gamma$-action to make $\varphi_{f}$ left
$\Gamma$-invariant.  The upshot is that we can study modular forms by
studying the spaces of functions that arise through the construction
\eqref{eq:corresp}.

Of course, not every $\varphi \colon \Gamma \backslash G \rightarrow
\C$ will arise as $\varphi_{f}$ for some $f\in S_{K} (\Gamma)$: after
all, $f$ is holomorphic and satisfies rather stringent growth
conditions.  Pinning down all the requirements is somewhat technical
and is (mostly) done in the sequel.

%\subsection{Arithmetic groups}\label{ss:a.gps}
\subsection{}\label{ss:a.gps}
%
Before we define automorphic forms, we need to find the correct
generalizations of our groups $\SL_{2} (\R)$ and $\Gamma_{0} (N)$.
The correct setup is rather technical, but this really reflects the
power of the general theory, which handles so many different
situations (e.g., Maass forms, Hilbert modular forms, Siegel
modular forms, etc.).

Let $G$ be a connected Lie group, and let $K\subset G$ be a maximal
compact subgroup.  We assume that $G$ is the set of real points of a
connected semisimple algebraic group $\bG $ defined over $\Q$.  These
conditions mean the following \cite[\S2.1.1]{plat.rapin}:
\begin{enumerate}
\item The group $\bG$ has the structure of an affine algebraic variety given by an ideal
$I$ in the ring $R = \C [x_{ij}, D^{-1} ]$, where the variables $\{
x_{ij} \mid 1\leq i,j\leq n\}$ should be interpreted as the entries of
an ``indeterminate matrix,'' and $D$ is the polynomial $\det
(x_{ij})$.  Both the group multiplication $\bG \times \bG \rightarrow \bG$
 and inversion $\bG \rightarrow \bG$ are required to be morphisms of
algebraic varieties.

The ring $R$ is the coordinate ring of the algebraic group $\GL_{n}$.
Hence this condition means that $\bG$ can be essentially viewed as a
subgroup of $\GL_{n} (\C)$ defined by polynomial equations in the
matrix entries of the latter.
\item \defn{Defined over $\Q$} means that $I$ is generated by
polynomials with rational coefficients.
\item \defn{Connected} means that $\bG $ is connected as an algebraic variety.
\item \defn{Set of real points} means that $G$ is the set of real
solutions to the equations determined by $I$. We write $G = \bG (\R)$.
\item \defn{Semisimple} means that the maximal connected 
solvable normal subgroup of $\bG$ is trivial.
\end{enumerate}

\begin{example}
The most important example for our purposes is 
the \defn{split
form of $\SL_{n}$}.  For this choice we have 
\[
\text{$G=\SL_{n} (\R)$ and $K=\SO (n)$.}
\]
\end{example}

\begin{example}\label{ex:resofscalars}
Let $F/\Q$ be a number field.  Then there is a $\Q$-group $\bG$ such
that $\bG (\Q) = \SL_{n} (F)$.  The group $\bG$ is constructed as
$\bR_{F/\Q } (\SL_{n})$, where $\bR_{F/\Q }$ denotes the
\defn{restriction of scalars} from $F$ to $\Q$
\cite[\S2.1.2]{plat.rapin}.  For example, if $F$ is totally real, the
group $\bR_{F/\Q} (\SL_{2})$ appears when one studies Hilbert modular
forms.

Let $(r,s)$ be the signature of the field $F$, so that $F\otimes \R
\simeq \R^{r}\times \C^s$.  Then $G = \SL_n (\R)^{r} \times \SL_n
(\C)^{s}$ and $K = \SO (n)^{r}\times \SU(n)^{s}$.
\end{example}


\begin{example}
Another important example is the \defn{split symplectic group}
$\SP_{2n}$.  This is the group that arises when one studies Siegel
modular forms.  The group of real points $\SP_{2n} (\R)$ is the
subgroup of $\SL_{2n} (\R)$ preserving a fixed nondegenerate
alternating bilinear form on $\R^{2n}$.  We have $K = \UN (n)$.
\end{example}

\subsection{}
%
To generalize $\Gamma_{0} (N)$, we need the notion of an
\defn{arithmetic group}.  This is a discrete subgroup $\Gamma$ of the
group of rational points $\bG (\Q)$ that is commensurable with the set
of integral points $\bG (\Z)$.  Here commensurable simply means that
$\Gamma \cap \bG (\Z)$ is a finite index subgroup of both $\Gamma$ and
$\bG (\Z)$; in particular $\bG (\Z)$ itself is an arithmetic group.

\begin{example}\label{ex:cong.subgp}
For the split form of $\SL_{n}$ we have $\bG (\Z) = \SL_{n} (\Z)
\subset \bG (\Q) = \SL_{n} (\Q)$.  A trivial way to obtain other
arithmetic groups is by conjugation: if $g\in \SL_{n} (\Q)$, then
$g\cdot \SL_{n} (\Z) \cdot g^{-1}$ is also arithmetic.

A more interesting collection of examples is given by the congruence
subgroups.  The \defn{principal congruence subgroup} $\Gamma (N)$ is
the group of matrices congruent to the identity modulo $N$ for some
fixed integer $N\geq 1$.  A \defn{congruence subgroup} is a group
containing $\Gamma (N)$ for some $N$.

In higher dimensions there are many candidates to generalize the Hecke
subgroup $\Gamma_{0} (N)$.  For example, one can take the
subgroup of $\SL_{n} (\Z)$ that is upper-triangular mod $N$.  From
a computational perspective, this choice is not so good since its
index in $\SL_{n} (\Z)$ is large.  A better choice, and the one that
usually appears in the literature, is to define $\Gamma_{0} (N)$ to be
the subgroup of $\SL_{n} (\Z)$ with bottom row congruent to $(0,\dotsc
,0,*) \mod N$.
\end{example}

%\subsection{Automorphic forms}\label{ss:aut.forms} We are finally
\subsection{}\label{ss:aut.forms} 
%
We are almost ready to define automorphic forms.  Let $\fg$ be the Lie
algebra of $G$, and let $U (\fg)$ be its universal enveloping algebra
over $\C$.  Geometrically, $\fg$ is just the tangent space at the
identity of the smooth manifold $G$.  The algebra $U (\fg)$ is a
certain complex associative algebra canonically built from $\fg$.  The
usual definition would lead us a bit far afield, so we will settle for an
equivalent characterization: $U (\fg)$ can be realized as a certain
subalgebra of the ring of differential operators on $C^{\infty} (G)$,
the space of smooth functions on $G$.

In particular, $G$ acts on $C^{\infty} (G)$ by \defn{left
translations}: given $g\in G$ and $f\in C^{\infty} (G)$, we define
\[
L_{g} (f) (x) := f (g^{-1}x).
\]
Then $U (\fg)$ can be identified with the ring of all differential
operators on $C^{\infty} (G)$ that are invariant under left
translation.  For our purposes the most important part of $U (\fg)$ is
its center $Z (\fg)$.  In
terms of differential operators, $Z (\fg)$ consists of those operators
that are also invariant under \defn{right translation}:
\[
R_{g} (f) (x) := f (xg).
\]

\begin{definition}\label{def:autform}
An \defn{automorphic form} on $G$ with
respect to $\Gamma$ is a function $\varphi \colon G\rightarrow \C$ satisfying
\begin{enumerate}
\item $\varphi (\gamma g) = \varphi (g)$ for all $\gamma \in \Gamma $,
\item the right translates $\{\varphi  (gk)\mid k\in K \}$ span a finite-dimensional
space $\xi$ of functions,\label{ktype}
\item there exists an ideal $J\subset Z (\fg)$ of finite codimension
such that $J$ annihilates $\varphi $, and 
\item $\varphi $ satisfies a certain growth condition that we do not
wish to make precise.  (In the literature, $\varphi $ is said to be
\defn{slowly increasing}.)
\end{enumerate}
\end{definition}

For fixed $\xi$ and $J$, we denote by $\sA(\Gamma ,\xi ,J,K)$ the
space of all functions
satisfying the above four conditions.  It is a basic theorem, due to
Harish-Chandra \cite{harish-chandra}, that $\sA(\Gamma ,\xi ,J,K)$ is
finite-dimensional.

\begin{example}\label{ex:holo.aut}
We can identify the cuspidal modular forms $S_{k} (N)$ in the language
of Definition \ref{def:autform}.  Given a modular form $f$, let
$\varphi_{f}\in C^{\infty} (\SL_{2} (\R))$ be the function from  
\eqref{eq:corresp}.  Then the map $f\mapsto \varphi_{f}$ identifies
$S_{k} (N)$ with the subspace $\sA_{k} (N)$ of functions $\varphi$ satisfying
\begin{enumerate}
\item $\varphi (\gamma g) = \varphi (g)$ for all $\gamma \in \Gamma_{0} (N)$,
\item $\varphi (gm) = \zeta_{m}^{-k}\varphi (g)$ for all $m\in \SO (2)$,
\item $(\Delta +\lambda_{k})\varphi = 0$, where $\Delta \in Z (\fg)$
is the \defn{Laplace--Beltrami--Casimir operator} and 
\[
\lambda_{k} = \frac{k}{2}\left(\frac{k}{2}-1 \right),
\]
\item $\varphi$ is slowly increasing, and 
\item $\varphi$ is \defn{cuspidal}.
\end{enumerate}

% the casimir operator is
% h^{2}/2 - h + 2ef
% it should be the same .. .. check it

The first four conditions parallel Definition \ref{def:autform}.  Item (1) is the
$\Gamma$-invariance.  Item (2) implies that the right translates of
$\varphi$ by $\SO (2)$ lie in a fixed finite-dimensional
representation of $\SO (2)$.  Item (3) is how holomorphicity appears,
namely that $\varphi$ is killed by a certain differential operator.
Finally, item (4) is the usual growth condition.

The only condition missing from the general
definition is (5), which is an extra constraint placed on $\varphi$ to
ensure that it comes from a cusp form.  This condition can be
expressed by the vanishing of certain integrals (``constant terms'');
for details we refer to \cite{bump.autforms, gelbart}.
\end{example}

\begin{example}
Another important example appears when we set $k=0$ in (2) in Example
\ref{ex:holo.aut} and relax (3) by requiring only that $(\Delta
-\lambda)\varphi =0$ for \emph{some} nonzero $\lambda \in \R$.  Such
automorphic forms cannot possibly arise from modular forms, since
there are no nontrivial cusp forms of weight 0.  However, there are
plenty of solutions to these conditions: they correspond to
\defn{real-analytic} cuspidal modular forms of weight 0 and
are known as \defn{Maass forms}.  Traditionally one writes $\lambda =
(1-s^{2})/4$.  The positivity of $\Delta$ implies that $s\in
(-1,1)$ or is purely imaginary.

Maass forms are highly elusive objects.  Selberg proved that there are
infinitely many linearly independent Maass forms of full level
(i.e., on $\SL_{2} (\Z)$), but to this date no explicit construction of
a single one is known.  (Selberg's argument is indirect and relies on
the trace formula; for an exposition see \cite{sarnak}.)  For higher
levels some explicit examples can be constructed using theta series
attached to indefinite quadratic forms \cite{vigneras}.  Numerically
Maass forms have been well studied; see for example \cite{farmer}.

In general the arithmetic nature of the eigenvalues
$\lambda$ that correspond to Maass forms is unknown, although a famous
conjecture of Selberg states that for congruence subgroups they
satisfy the inequality $\lambda \geq 1/4$ (in other words, only purely
imaginary $s$ appear above).  The truth of this
conjecture would have far-reaching consequences, from analytic number
theory to graph theory \cite{lubotzky}.
\end{example}

\subsection{}\label{ss:cusp.eis}
As Example \ref{ex:holo.aut} indicates, there is a notion of \defn{cuspidal
automorphic form}.  The exact definition is too technical to state
here, but it involves an appropriate generalization of the notion of
constant term familiar from modular forms.

There are also \defn{Eisenstein series} \cite{langlands.e.s,
arthur.e.s}.  Again the complete definition is technical; we only
mention that there are different types of Eisenstein series
corresponding to certain subgroups of $G$.  The Eisenstein series that
are easiest to understand are those built from cusp forms on lower
rank groups.  Very explicit formulas for Eisenstein series on
$\GL_{3}$ can be seen in \cite{bump.gl3}.  For a down-to-earth
exposition of some of the Eisenstein series on $\GL_{n}$, we refer to
\cite{goldfeld}.

The decomposition of $M_{k} (\Gamma_{0} (N))$ into cusp forms and
Eisenstein series also generalizes to a general group $G$, although
the statement is much more complicated.  The result is a theorem of
Langlands \cite{langlands} known as the {\em spectral decomposition
of $L^{2} (\Gamma \backslash G)$}.  A thorough recent presentation of
this can be found in \cite{mw}.

%\subsection{Cohomology of arithmetic groups}\label{ss:coho}
\subsection{}\label{ss:coho}
%
Let $\sA = \sA (\Gamma ,K)$ be the space of all automorphic forms,
where $\xi$ and $J$ range over all possibilities.  The space $\sA$ is
huge, and the arithmetic significance of much of it is unknown.  This
is already apparent for $G=\SL_{2} (\R)$.  The automorphic forms
directly connected with arithmetic are the holomorphic modular forms,
not the Maass forms\footnote{However, Maass forms play a very
important \emph{indirect} role in arithmetic.}. Thus the question
arises: which automorphic forms in $\sA$ are the most natural
generalization of the modular forms?

One answer is provided by the isomorphisms \eqref{eq:es},
\eqref{eq:eshw}, \eqref{eq:eseis}.  These show that modular forms
appear naturally in the cohomology of modular curves.  Hence a
reasonable approach is to generalize the \emph{left} of \eqref{eq:es},
\eqref{eq:eshw}, \eqref{eq:eseis}, and to study the resulting
cohomology groups.  This is the approach we will take.  One drawback
is that it is not obvious that our generalization has anything to do
with automorphic forms, but we will see eventually that it certainly
does.  So we begin by looking for an appropriate generalization of the
modular curve $Y_{0} (N)$.

Let $G$ and $K$ be as in Section~\ref{ss:a.gps}, and let $X$ be the quotient
$G/K$.  This is a global Riemannian symmetric space \cite{helgason}.  One can prove
that $X$ is contractible.  Any arithmetic group $\Gamma\subset G$ acts
on $X$ properly discontinuously.  In particular, if $\Gamma$ is
torsion-free, then the quotient $\Gamma \backslash X$ is a smooth
manifold.

Unlike the modular curves, $\Gamma \backslash X$ will not have a
complex structure in general\footnote{The symmetric spaces that have
a complex structure are known as \defn{bounded domains}, or
\defn{Hermitian symmetric spaces} \cite{helgason}.}; nevertheless,
$\Gamma \backslash X$ is a very nice space.  In particular, if
$\Gamma$ is torsion-free, it is an \defn{Eilenberg--Mac~Lane} space for
$\Gamma$, otherwise known as a $K (\Gamma , 1)$.  This means that the
only nontrivial homotopy group of $\Gamma \backslash X$ is its
fundamental group, which is isomorphic to $\Gamma$, and that the
universal cover of $\Gamma \backslash X$ is contractible.  Hence
$\Gamma \backslash X$ is in some sense a ``topological 
incarnation''\footnote{This apt phrase is due to Vogan \cite{vogan}.}
of~$\Gamma$.

This leads us to the notion of the \defn{group cohomology} $H^{*}
(\Gamma ; \C)$ of $\Gamma$ with trivial complex coefficients.  In the
early days of algebraic topology, this was defined to be the complex
cohomology of an Eilenberg--Mac~Lane space for $\Gamma$
\cite[Introduction, I.4]{brown}:
\begin{equation}\label{eq:coho.iso}
H^{*} (\Gamma ;\C) = H^{*} (\Gamma \backslash X; \C).
\end{equation}
Today there are purely algebraic approaches to $H^{*} (\Gamma ;\C)$
\cite[III.1]{brown}, but for our purposes \eqref{eq:coho.iso} is
exactly what we need.  In fact, the group cohomology $H^{*} (\Gamma ;
\C)$ can be identified with the cohomology of the quotient $\Gamma
\backslash X$ even if $\Gamma$ has torsion, since we are working with
complex coefficients.  The cohomology groups $H^{*} (\Gamma ; \C)$,
where $\Gamma$ is an arithmetic group, are our proposed generalization
for the weight 2 modular forms.

What about higher weights?  For this we must replace the trivial
coefficient module $\C$ with local systems, just as we did in
\eqref{eq:eshw}.  For our purposes it is enough to let $\MMM$ be a
rational finite-dimensional representation of $G$ over the complex
numbers.  Any such $\MMM$ gives a representation of $\Gamma \subset
G$ and thus induces a local system $\widetilde{\MMM}$ on $\Gamma
\backslash X$.  As before, the group cohomology $H^{*} (\Gamma ;
\MMM)$ is the cohomology $H^{*} (\Gamma\backslash X;
\widetilde{\MMM})$.  In \eqref{eq:eshw} we took $\MMM =M_{n}$, the
$n$th symmetric power of the standard representation.  For a general
group $G$ there are many kinds of representations to consider.  In any
case, we contend that the cohomology spaces
\[
H^{*} (\Gamma ; \MMM) =  H^{*}(\Gamma \backslash X; \widetilde{\MMM}) 
\]
are a good generalization of the spaces of modular forms.

\subsection{}\label{ss:franke}
%
It is certainly not obvious that the cohomology groups $H^{*} (\Gamma
; \MMM)$ have \emph{anything} to do with automorphic forms, although
the isomorphisms \eqref{eq:es}, \eqref{eq:eshw}, \eqref{eq:eseis} look
promising.  

The connection is provided by a deep theorem of Franke \cite{franke},
which asserts
that
\begin{enumerate}
\item the cohomology groups $H^{*} (\Gamma ; \MMM)$ can be directly
computed in terms of certain automorphic forms (the automorphic forms
of ``cohomological type,'' also known as those with ``nonvanishing
$(\fg , K)$ cohomology'' \cite{vog.zuck}); and \label{bc}
\item there is a direct sum decomposition 
\begin{equation}\label{franke.decomp}
H^{*} (\Gamma ; \MMM) = H^{*}_{\text{cusp}} (\Gamma ;
\MMM)\oplus\bigoplus_{\{P \}}H^{*}_{\{P \}} (\Gamma ; \MMM),
\end{equation}
where the sum is taken over the set of classes of \defn{associate proper
$\Q$-parabolic subgroups of $G$}. \label{ld}
\end{enumerate}
The precise version of statement \eqref{bc} is known in the literature
as the \defn{Borel conjecture}.  Statement \eqref{ld} parallels
Langlands's spectral decomposition of $L^{2} (\Gamma \backslash G)$.

\begin{example}
For $\Gamma = \Gamma_{0} (N)\subset \SL_{2} (\Z)$, the decomposition
\eqref{franke.decomp} is exactly \eqref{eq:eseis}.  The cusp forms
$S_{k} (N)\oplus \overline{S_{k} (N)}$ correspond to the summand
$H^{1}_{\text{cusp}} (\Gamma ; \MMM)$.  There is one class of proper
$\Q$-parabolic subgroups in $\SL_{2} (\R)$, represented by the Borel subgroup
of upper-triangular matrices.  Hence only one term appears in big
direct sum on the right of \eqref{franke.decomp}, which is the
Eisenstein term
$\Eis_{k}$.
%For general $G$,  \eqref{franke.decomp}
%parallels the Hence for these $\MMM$ the cohomology $H^{*} (\Gamma ;
%\MMM)$ can be completely described in terms of automorphic forms.
\end{example}


The summand $H^{*}_{\text{cusp}} (\Gamma ; \MMM)$ of
\eqref{franke.decomp} is called the \defn{cuspidal cohomology}; this
is the subspace of classes represented by cuspidal automorphic forms.
The remaining summands constitute the \defn{Eisenstein cohomology} of
$\Gamma$ \cite{harder.icm}.  In particular the summand indexed by $\{P \}$ is
constructed using Eisenstein series attached to certain cuspidal
automorphic forms on lower rank groups.  Hence $H^{*}_{\text{cusp}}
(\Gamma ; \MMM)$ is in some sense the most important part of the
cohomology: all the rest can be built systematically from cuspidal
cohomology on lower rank groups\footnote{This is a bit of an
oversimplification, since it is a highly nontrivial problem to decide
when cusp cohomology from lower rank groups appears in $\Gamma$.
However, many results are known; as a selection we mention
\cite{harder.icm, harder.gl2, li.schwermer}}. This leads us to our
basic computational problem:

\begin{problem}
Develop tools to compute explicitly the cohomology spaces $H^{*}
(\Gamma ; \MMM)$ and to identify the cuspidal
subspace $H^{*}_{\text{cusp}} (\Gamma ; {\MMM})$. 
\end{problem}

%%%%%%%%%%%%%
%           %
%  retract  %
%           %
%%%%%%%%%%%%%

\section{Combinatorial Models for Group Cohomolo\-gy}\label{comb.models}
%
\subsection{}
%
In this section, we restrict attention to $G=\SL_{n} (\R)$ and
$\Gamma$, a congruence subgroup of $\SL_{n} (\Z)$.  By the previous
section, we can study the group cohomology $H^{*} (\Gamma ; \MMM)$ by
studying the cohomology $H^{*} (\Gamma \backslash X;
\widetilde{\MMM})$.  The latter spaces can be studied using standard
topological techniques, such as taking the cohomology of complexes
associated to cellular decompositions of $\Gamma \backslash X$.  For
$\SL_{n} (\R ) $, one can construct such decompositions using a
version of explicit reduction theory of real positive-definite
quadratic forms due to \Vor\ \cite{vor1}.  The goal of this section is
to explain how this is done.  We also discuss how the cohomology can
be explicitly studied for congruence subgroups of $\SL_{3} (\Z)$.

%\subsection{A linear model for the symmetric
%space}\label{vcd.section}
\subsection{}\label{vcd.section}
%
Let $V$ be the $\R $-vector space of all symmetric $n\times n$
matrices, and let $C\subset V$ be the subset of positive-definite
matrices.  The space $C$ can be identified with the space of all real
positive-definite quadratic forms in $n$ variables: in coordinates, if $x=
(x_{1},\dotsc ,x_{n})^{t}\in \R^{n}$ (column vector), then the matrix
$A\in C$ induces the 
quadratic form  
\[
x \longmapsto x^{t}Ax,
\]
and it is well known that any positive-definite quadratic form arises
in this way.  The space $C$ is a cone, in that it is preserved by
homotheties: if $x\in C$, then $\lambda x\in C$ for all $\lambda \in
\R_{>0}$.  It is also convex: if $x_{1}, x_{2}\in C$, then $tx_{1}+
(1-t)x_{2}\in C$ for $t\in [0,1]$.  Let $D$ be the quotient of $C$ by
homotheties.

\begin{example}
The case $n=2$ is illustrative.  We can take
coordinates on $V \simeq \R^{3}$ by representing any matrix in $V$ as 
\[
\left(\begin{array}{cc}
x&y\\
y&z
\end{array} \right), \quad x,y,z\in \R.  
\] 
The subset of singular matrices $Q = \{xz-y^{2}=0 \}$ is a quadric
cone in $V$ dividing the complement $V\smallsetminus Q$ into three
connected components.  The component containing the identity matrix is
the cone $C$ of positive-definite matrices. The quotient $D$ can be
identified with an open $2$-disk.
\end{example}

The group $G$ acts on $C$ on the left by
\[
(g,c)\longmapsto gcg^{t}.
\]
This action commutes with that of the homotheties and thus
descends to a $G$-action on $D$.  One can show that $G$ acts transitively
on $D$ and that the stabilizer of the image of the identity matrix is
$K=\SO (n)$.  Hence we may identify $D$ with our symmetric space $X =
\SL_{n} (\R )/\SO (n)$.  We will do this in the sequel, using the
notation $D$ when we want to emphasize the coordinates coming from the
linear structure of $C\subset V$ and using the notation $X$ for the
quotient $G/K$.

We can make the identification $D\simeq X$ more explicit.  If $g\in
\SL_{n} (\R)$, then the map 
\begin{equation}\label{chmap}
g \longmapsto gg^{t}
\end{equation}
takes $g$ to a symmetric positive-definite matrix.  Any coset $gK$ is
taken to the same matrix since $KK^{t} = \Id$.  Thus \eqref{chmap}
identifies $G/K$ with a subset $C_{1}$ of $C$, namely those
positive-definite symmetric matrices with determinant $1$.  It is
easy to see that $C_{1}$ maps diffeomorphically onto $D$.  

The inverse map $C_{1}\rightarrow G/K$ is more complicated. Given a
determinant $1$ positive-definite symmetric matrix $A$, one must find
$g\in \SL_{n} (\R)$ such that $gg^{t}=A$.  Such a representation
always exists, with $g$ determined uniquely up to right multiplication
by an element of $K$.  In computational linear algebra, such a $g$ can
be constructed through \defn{Cholesky decomposition} of $A$.

The group $\SL_{n} (\Z )$ acts on $C$ via the $G$-action and does so
properly discontinuously.  This is the ``unimodular change of
variables'' action on quadratic forms \cite[V.1.1]{serre:arithmetic}.
Under our identification of $D$ with $X$, this is the usual action of
$\SL_{n} (\Z)$ by left translation from Section~\ref{ss:coho}.

\subsection{}\label{ss:vcd}
%
Now consider the group cohomology $H^{*} (\Gamma ; \MMM) = H^{*}
(\Gamma \backslash X; \widetilde{\MMM})$.  The identification $D\simeq
X$ shows that the dimension of $X$ is $n (n+1)/2 - 1$.  Hence $H^{i}
(\Gamma ; \MMM)$ vanishes if $i > n (n+1)/2 - 1$.  Since $\dim X$
grows quadratically in $n$, there are many potentially interesting
cohomology groups to study.

However, it turns out that there is some additional vanishing of the
cohomology for deeper (topological) reasons.  For $n=2$, this is easy
to see.  The quotient $\Gamma \backslash \fh$ is homeomorphic to a
topological surface with punctures, corresponding to the cusps of
$\Gamma$.  Any such surface $S$ can be retracted onto a finite graph
simply by ``stretching'' $S$ along its punctures.  Thus $H^{2} (\Gamma
; \MMM) = 0$, even though $\dim \Gamma \backslash \fh = 2$.

For $\Gamma \subset \SL_{n} (\Z )$, a theorem of Borel--Serre implies
that $H^{i} (\Gamma ;{\MMM})$ vanishes if $i> \dim X -n+1 = n(n-1)/2$
\cite[Theorem 11.4.4]{borel.serre}.  The number $\nu = n (n-1)/2$ is
called the \defn{virtual cohomological dimension} of $\Gamma $ and is
denoted $\vcd \Gamma$.  Thus we only need to consider cohomology in
degrees $i\leq \nu$.

Moreover we know from Section~\ref{ss:franke} that the most interesting part
of the cohomology is the cuspidal cohomology.  In what degrees can it
live?  For $n=2$, there is only one interesting cohomology group
$H^{1} (\Gamma ; \MMM)$, and it contains the cuspidal cohomology.  For
higher dimensions, the situation is quite different: for most $i$, the
subspace $H^{i}_{\text{cusp}} (\Gamma ; \MMM)$ vanishes!  In fact in
the late 1970's Borel, Wallach, and Zuckerman observed that the cuspidal
cohomology can only live in the cohomological degrees lying in an
interval around $(\dim X)/2$ of size linear in $n$.  An explicit
description of this interval is given in \cite[Proposition
3.5]{schwermer}; one can also look at Table \ref{dimtable}, from which
the precise statement is easy to determine.

Another feature of Table \ref{dimtable} deserves to be mentioned.
There are exactly two values of $n$, namely $n=2, 3$, such that
virtual cohomological dimension equals the upper limit of the cuspidal
range.  This will have implications later, when we study the action of
the Hecke operators on the cohomology.

\begin{table}[htb]
\begin{center}
\begin{tabular}{c||c|c|c|c|c|c|c|c}
$n$&2&3&4&5&6&7&8&9\cr
\hline
\hline
$\dim X$&2 & 5 & 9 & 14 & 20 & 27 & 35 & 44 \cr
$\vcd \Gamma $ & 1 & 3 & 6 & 10 & 15 & 21 & 28 & 36\cr
top degree of $H^{*}_{\text{cusp}} $& 1 & 3 & 5 & 8 & 11 & 15 & 19& 24 \cr
bottom degree of $H^{*}_{\text{cusp}} $& 1 & 2 & 4 & 6 & 9 & 12 & 16 & 20 \cr
\end{tabular}
\end{center}
\medskip
\caption{The virtual cohomological dimension and the cuspidal range
for subgroups of $\SL_{n}(\Z)$.\label{dimtable}}
\end{table}

%\subsection{The \Vor\  polyhedron}\label{vor.poly.section}
\subsection{}\label{vor.poly.section}
%
Recall that a point in $\Z ^{n}$ is said to be \defn{primitive} if the
greatest common divisor of its coordinates is $1$.  In particular, a
primitive point is nonzero.  Let $\PPP\subset \Z ^{n}$ be the set of
primitive points.  Any $v\in \PPP$, written as a column vector,
determines a rank-$1$ symmetric matrix $q (v)$ in the closure $\bar C$
via $q ( v) = vv^{t}$.  The \defn{\Vor \ polyhedron} $\Pi $ is defined
to be the closed convex hull in $\bar C$ of the points $q (v)$, as $v$
ranges over $\PPP$.  Note that by construction, $\SL_{n} (\Z )$ acts
on $\Pi $, since $\SL_{n} (\Z)$ preserves the set $\{q (v) \}$ and
acts linearly on $V$.

\begin{example}\label{ex:vor.poly}
Figure \ref{tessfig} represents a crude attempt to show what $\Pi$
looks like for $n=2$.  These images were constructed by computing a
large subset of the points $q (v)$ and taking the convex hull (we took
all points $v\in \PPP$ such that $\Trace q (v) < N$ for some large
integer $N$).  From a distance, the polyhedron $\Pi$ looks almost
indistinguishable from the cone $C$; this is somewhat conveyed by the
right of Figure \ref{tessfig}.  Unfortunately $\Pi$ is not locally
finite, so we really cannot produce an accurate picture.  To get a more
accurate image, the reader should imagine that each vertex meets
infinitely many edges.  On the other hand, $\Pi$ is not hopelessly
complex: each maximal face is a triangle, as the pictures suggest.


\begin{figure}[htb] \centering
\subfigure[\label{fig.origin}]{\includegraphics[scale=0.2]{diagrams/tess1}}
\quad\quad
\subfigure[\label{fig.side}]{\includegraphics[scale=0.2]{diagrams/tess2}}
\caption{\label{tessfig}The polyhedron $\Pi$ for $\SL_{2} (\Z)$.  In~(a) we see $\Pi$ from the origin, in~(b) from
the side.  The small triangle at the right center of~(a) is the
facet with vertices $\{q (e_{1}), q (e_{2}), q (e_{1}+e_{2}) \}$,
where $\{e_{1},e_{2}\}$ is the standard basis of $\Z^{2}$.  In~(b) the $x$-axis runs along the top from left to right, and
the $z$-axis runs down the left side.  The facet from~(a) is the little triangle at the top left corner of~(b).}
\end{figure}
\end{example}


\subsection{}
%
The polyhedron $\Pi$ is quite complicated: it has infinitely many
faces and is not locally finite.  However, one of \Vor 's great
insights is that $\Pi$ is actually not as complicated as it seems.

For any $A\in C$, let $\mu (A)$ be the minimum value attained by $A$
on $\PPP$ and let $M (A)\subset \PPP$ be the set on which $A$ attains
$\mu (A)$.  Note that $\mu (A)> 0 $ and $M (A)$ is finite since $A$ is
positive-definite.  Then $A$ is called \defn{perfect} if it is
recoverable from the knowledge of the pair $(\mu (A), M (A) )$.  In
other words, given $(\mu (A), M (A) )$, we can write a system of
linear equations
\begin{equation}\label{system}
mZm^{t} = \mu (A), \quad m\in M (A),
\end{equation}
where $Z = (z_{ij})$ is a symmetric matrix of variables.  Then $A$ is
perfect if and only if $A$ is the unique solution to the system
\eqref{system}. 

\begin{example}\label{ex:A2}
The quadratic form $Q (x,y) =  x^{2}-xy+y^{2}$ is perfect.  The smallest
nontrivial value
it attains on $\Z^{2}$ is $\mu (Q) = 1$, and it does so on the columns of 
\[
M (Q) =\left(\begin{array}{ccc}
1&0&1\\
0&1&1
\end{array} \right)
\]   
and their negatives.   Letting $\alpha x^{2}+\beta xy + \gamma
y^{2}$ be an undetermined quadratic form and applying the data $(\mu
(Q), M (Q))$, we are led to the system of linear equations 
\[
\alpha =1, \quad \gamma =1, \quad \alpha +\beta +\gamma =1.
\]
From this we recover $Q (x,y)$.
\end{example}

\begin{example}\label{ex:cubicform}
The quadratic form $Q' (x,y) = x^{2}+y^{2}$ is not perfect.  Again the
smallest nontrivial value of $Q'$ on $\Z^{2}$ is $m (Q')= 1$, attained on the
columns of 
\[
M (Q')=\left (\begin{array}{cc}
1&0\\
0&1
\end{array}\right)
\]
and their negatives.  But every member of the one-parameter family of quadratic forms 
\begin{equation}\label{}
x^{2}+\alpha xy+y^{2}, \quad \alpha \in (-1,1)
\end{equation}
has the same
set of minimal vectors, and so $Q'$ cannot be recovered from the
knowledge of $m (Q'), M (Q')$.
\end{example}

\begin{example}\label{ex:An}
Example \ref{ex:A2} generalizes to all $n$.  Define
\begin{equation}\label{anform}
A_{n} (x) := \sum_{i=1}^{n}x_{i}^{2} - \sum_{1\leq i<j\leq n} x_{i}x_{j}.
\end{equation}
Then $A_{n}$ is perfect for all $n$.  We have $\mu
(A_{n})=1$, and $M (A_{n})$ consists of all points of the form
\[
\pm (e_{i}+e_{i+1}+\dotsb + e_{i+k}), \quad 1\leq i\leq n, \quad i\leq i+k\leq n,
\]
where $\{e_{i} \}$ is the standard basis of $\Z^{n}$.  This quadratic
form is closely related to the $A_{n}$ root lattice
\cite{fulton.harris}, which explains its name.  It is one of two
infinite families of perfect forms studied by \Vor (the other is
related to the $D_{n}$ root lattice).
\end{example}

We can now summarize \Vor 's main results:
\begin{enumerate}
\item There are finitely many equivalence classes of perfect forms
modulo the action of $\SL_{n} (\Z)$.  \Vor \ even gave an explicit
algorithm to determine all the perfect forms of a given dimension.
\item \label{item:minvec} The facets of $\Pi$, in other words the
codimension $1$ faces, are in bijection with the rank $n$ perfect
quadratic forms.  Under this correspondence the minimal vectors $M
(A)$ determine a facet $F_{A}$ by taking the convex hull in $\bar C$
of the finite point set $\{q (m)\mid m\in M (A) \}$.  Hence there are
finitely many faces of $\Pi$ modulo $\SL_{n} (\Z)$ and thus finitely
many modulo any finite index subgroup $\Gamma$.
\item Let $\V$ be the set of cones over the faces of $\Pi $.  Then $\V$
is a \defn{fan}, which means (i) if $\sigma \in \V$, then any face of
$\sigma$ is also in $\V$; and (ii) if $\sigma ,\sigma '\in \V$, then
$\sigma \cap \sigma '$ is a common face of each\footnote{Strictly
speaking, \Vor \ actually showed that every codimension 1 cone is
contained in two top-dimensional cones.}.  The fan $\V$ provides a
reduction theory for $C$ in the following sense: any point $x\in C$ is
contained in a unique $\sigma (x) \in \V$, and the set $\{\gamma \in
\SL _{n} (\Z ) \mid \gamma \cdot \sigma (x) = \sigma (x) \}$ is
finite.  \Vor\ also gave an explicit algorithm to determine $\sigma
(x)$ given $x$, the \defn{\Vor \ reduction algorithm}.
\end{enumerate}

The number $N_{\text{perf}}$ of equivalence classes of perfect forms modulo the
action of $\GL_{n}(\Z)$ grows rapidly with
$n$ (Table \ref{vortable}); the complete classification is known
only for $n\leq 8$.  For a list of perfect forms up to $n = 7$, see
\cite{cs}.  For a recent comprehensive treatment of perfect forms,
with many historical remarks, see
\cite{martinet}.

%\begin{table}[htb]
%\begin{center}
%\begin{tabular}{c|ccccccc}
%$n$&2&3&4&5&6&7&8\\
%\hline
%$N$&1&1&2&3&?&33&10916\\
%\end{tabular}
%\end{center}
%\medskip
%\caption{\label{vortable}}
%\end{table}

\begin{table}[htb]
\begin{center}
\begin{tabular}{c|c|c}
Dimension&$N_{\text{perf}}$&Authors\\
\hline
\hline
$2$&$1$&\Vor \cite{vor1}\\
$3$&$1$& \emph{ibid.}\\
$4$&$2$& \emph{ibid.}\\
$5$&$3$& \emph{ibid.}\\
$6$&$7$&Barnes \cite{barnes}\\
$7$&$33$&Jaquet-Chiffelle \cite{jaquet, jaquet-chif}\\
$8$&$10916$&Dutour--Sch\"urmann--Vallentin \cite{dutour}\\
\end{tabular}
\end{center}
\medskip
\caption{The number $N_{\text{perf}}$ of equivalence 
classes of perfect forms.\label{vortable}}
\end{table}


%\subsection{Cellular decompositions and the well-rounded
%retract}\label{wrr.section}
\subsection{}\label{wrr.section}
%
Our goal now is to describe how the \Vor \
fan $\V$ can be used to compute the cohomology $H^{*} (\Gamma ;
\MMM)$.  The idea is to use the cones in $\V$ to chop the quotient $D
$ into pieces.  

For any $\sigma \in \V$, let $\sigma^{\circ}$ be the open cone
obtained by taking the complement in $\sigma$ of its proper faces.
Then after taking the quotient by homotheties, the cones $\{\sigma^{\circ }
\cap C \mid \sigma \in \V \}$ pass to locally closed subsets of $D$.
Let $\sC$ be the set of these images.

Any $c\in \sC$ is a \defn{topological cell}, i.e., it is homeomorphic to
an open ball, since $c$ is homeomorphic to a face of $\Pi$.  Because
$\sC$ comes from the fan $\V$, the cells in $\sC$ have good incidence
properties: the closure in $D$ of any $c\in \sC$ can be written as a
finite disjoint union of elements of $\sC$.  Moreover, $\sC$ is
locally finite: by taking quotients of all the $\sigma^{\circ }$
meeting $C$, we have eliminated the open cones lying in $\bar C$, and
it is these cones that are responsible for the failure of local
finiteness of $\V$.  We summarize these properties by saying that
$\sC$ gives a \defn{cellular decomposition} of $D$.  Clearly $\SL_{n}
(\Z)$ acts on $\sC$, since $\sC$ is constructed using the fan $\V$.
Thus we obtain a cellular decomposition\footnote{If $\Gamma$ has torsion, then
cells in $\sC$ can have nontrivial stabilizers in $\Gamma$, and thus
$\Gamma \backslash \sC$ should be considered as an ``orbifold''
cellular decomposition.} of $\Gamma \backslash D $ for
any torsion-free $\Gamma$.  We call $\sC$ the \defn{\Vor \
decomposition} of $D$.

Some care must be taken in using these cells to perform topological
computations.  The problem is that even though the individual pieces
are homeomorphic to balls and are glued together nicely, the
boundaries of the closures of the pieces are not homeomorphic to
spheres in general.  (If they were, then the \Vor \ decomposition
would give rise to a \defn{regular} cell complex \cite{cooke.finney},
which can be used as a substitute for a simplicial or CW complex in
homology computations.)  Nevertheless, there is a way to remedy this.

Recall that a subspace $A$ of a topological space $B$ is a
\defn{strong deformation retract} if there is a continuous map
$f\colon B\times [0,1]\rightarrow B$ such that $f (b,0)=b$, $f
(b,1)\in A$, and $f (a,t) = a$ for all $a\in A$.  For such pairs
$A\subset B$ we have $H^{*} (A) = H^{*} (B)$.  One can show that there
is a strong deformation retraction from $C$ to itself equivariant
under the actions of both $\SL_n(\Z)$ and the homotheties and that
the image of the retraction modulo homotheties, denoted $W$, is
naturally a locally finite regular cell complex of dimension $\nu $.
Moreover, the cells in $W$ are in bijective, inclusion-reversing
correspondence with the cells in $\sC $.  In particular, if a cell in
$\sC$ has \defn{codimension} $d$, the corresponding cell in $W$ has
\defn{dimension} $d$.  Thus, for example, the vertices of $W$ modulo
$\SL_{n} (\Z)$ are in bijection with the top-dimensional cells in
$\sC$, which are in bijection with equivalence classes of perfect
forms.

In the literature $W$ is called the \defn{well-rounded retract}.  The
subspace $W\subset D \simeq X$ has a beautiful geometric
interpretation.  The quotient
\[
\SL_{n} (\Z)\backslash X = \SL_{n} (\Z)
\backslash \SL_{n} (\R)/ \SO (n)
\]
can be interpreted as the moduli space of lattices in $\R^{n}$ modulo
the equivalence relation of rotation and positive scaling
(cf. \cite{ash.gross}; for $n=2$ one can also see \cite[VII,
Proposition 3]{serre:arithmetic}).  Then $W$ corresponds to those lattices whose
shortest nonzero vectors span $\R^{n}$.  This is the origin of the
name: the shortest vectors of such a lattice are ``more round'' than
those of a generic lattice.

The space $W$ was known classically for $n=2$ and was constructed for
$n\geq 3$ by Lannes and Soul\'e, although Soul\'e only published the
case $n=3$ \cite{soule}.  The construction for all $n$ appears in work
of Ash \cite{Ash80, ash.wrr}, who also generalized $W$ to a much
larger class of groups.  Explicit computations of the cell structure
of $W$ have only been performed up to $n=6$ \cite{gangl}.  Certainly
computing $W$ explicitly for $n=8$ seems very difficult, as Table
\ref{vortable} indicates.

\begin{example}\label{sl2}
Figure \ref{tess} illustrates $\sC$ and $W$ for $\SL_{2} (\Z)$.  As in
Example \ref{ex:vor.poly}, the polyhedron $\Pi$ is 3-dimensional, and
so the \Vor \ fan $\V$ has cones of dimensions $0,1,2,3$.  The
$1$-cones of $\V$, which correspond to the vertices of $\Pi$, pass to
infinitely many points on the boundary $\partial \bar D = \bar D
\smallsetminus D$.  The $3$-cones become triangles in $\bar D$ with
vertices on $\partial \bar D$.  In fact, the identifications $D\simeq
\SL_{2} (\R)/\SO (2)\simeq \fh$ realize $D$ as the Klein model for the
hyperbolic plane, in which geodesics are represented by Euclidean line
segments.  Hence, the images of the $1$-cones of $\V$ are none other
than the usual cusps of $\fh$, and the triangles are the $\SL_{2}
(\Z)$-translates of the ideal triangle with vertices $\{0,1,\infty
\}$.  These triangles form a tessellation of $\fh$ sometimes known as
the \defn{Farey tessellation}. The edges of the \Vor are the $\SL_{2}
(\Z)$-translates of the ideal geodesic between $0$ and $\infty$.
After adjoining cusps and passing to the quotient $X_{0} (N)$, these
edges become the supports of the Manin symbols from Section~\ref{sec:manin}
(cf.~Figure \ref{figure:modsym}).  This example also shows how the
\Vor \ decomposition fails to be a regular cell complex: the
boundaries of the closures of the triangles in $D$ do not contain the
vertices and thus are not homeomorphic to circles.

The virtual cohomological dimension of $\SL_{2} (\Z)$ is 1.  Hence the
well-rounded retract $W$ is a graph (Figures \ref{tess} and
\ref{retract2.fig}).  Note that $W$ is not a manifold.  The vertices
of $W$ are in bijection with the Farey triangles---each vertex lies at
the center of the corresponding triangle---and the edges are in
bijection with the Manin symbols.  Under the map $D\rightarrow \fh$,
the graph $W$ becomes the familiar ``$\PSL_{2}$-tree'' embedded in
$\fh$, with vertices at the order 3 elliptic points (Figure
\ref{retract2.fig}).

\begin{figure}[htb]
\begin{center}
\includegraphics[scale=0.4]{diagrams/retract3}
\end{center}
\caption{\label{tess}The \Vor \ decomposition and the retract in $D$.}
\end{figure}

\begin{figure}[htb]
\begin{center}
\includegraphics[scale=0.4, trim=0 330 0 0]{diagrams/retract2}
\end{center}
\caption{\label{retract2.fig}The \Vor \ decomposition and the retract in $\fh$.}
\end{figure}

\end{example}

%\subsection{Example: $\SL_{3} (\Z)$}\label{sl3}
\subsection{}\label{sl3}

We now discuss the example $\SL_{3} (\Z)$ in some detail.  This
example gives a good feeling for how the general situation compares to
the case $n=2$.

We begin with the \Vor \ fan $\V$.  The cone $C$ is 6-dimensional, and
the quotient $D$ is 5-dimensional.  There is one equivalence class of
perfect forms modulo the action of $\SL_{3} (\Z )$, represented by the
form \eqref{anform}.  Hence there are 12 minimal vectors; six are the
columns of the matrix
\begin{equation}\label{sl3minvecs}
\left(
\begin{array}{cccccc}
1&0&0&1&0&1\\
0&1&0&1&1&1\\
0&0&1&0&1&1
\end{array}
 \right),
\end{equation}
and the remaining six are the negatives of these.  This implies that
the cone $\sigma$ corresponding to this form is 6-dimensional and
simplicial.  The latter implies that the faces of $\sigma$ are the
cones generated by $\{q (v)\mid v\in S \}$, where $S$ ranges over all
subsets of \eqref{sl3minvecs}.  To get the full structure of the fan,
one must determine the $\SL_{3} (\Z)$ orbits of faces, as well as
which faces lie in the boundary $\partial \bar C = \bar C
\smallsetminus C$.  After some pleasant computation, one finds:
\begin{enumerate}
\item There is one equivalence class modulo $\SL_{3} (\Z)$ for each of the
6-, 5-, 2-, and
1-dimensional cones.
\item There are two equivalence classes of the 4-dimensional cones,
represented by the sets of minimal vectors
\[
\left(\begin{array}{cccc}
1&0&0&1\\
0&1&0&1\\
0&0&1&1
\end{array} \right)\quad\text{and}\quad 
\left(\begin{array}{cccc}
1&0&0&1\\
0&1&0&1\\
0&0&1&0
\end{array} \right).\label{3faces}
\]
\item There are two equivalence classes of the 3-dimensional cones,
represented by the sets of minimal vectors
\[
\left(\begin{array}{cccc}
1&0&0\\
0&1&0\\
0&0&1
\end{array} \right)\quad\text{and}\quad 
\left(\begin{array}{cccc}
1&0&1\\
0&1&1\\
0&0&0
\end{array} \right).
\]
The second type of 3-cone lies in $\partial \bar C$ and thus does 
not determine a cell in $\sC$.

\item The 2- and 1-dimensional cones lie entirely in
$\partial \bar C$ and do not determine cells in $\sC$.
\end{enumerate}

After passing from $C$ to $D$, the cones of dimension $k$ determine
cells of dimension $k-1$.
Therefore, modulo the action of $\SL_{3} (\Z)$ there are five types of
cells in the \Vor \ decomposition $\sC$, with dimensions from $5$ to $2$.
We denote these cell types by $c_{5}$, $c_{4}$, $c_{3a}$, $c_{3b}$,
and $c_{2}$.  Here $c_{3a}$ corresponds to the first type of 4-cone in item
\eqref{3faces} above, and $c_{3b}$ to the second.  For a beautiful way
to index the cells of $\sC$ using configurations in projective spaces,
see \cite{M91}.

The virtual cohomological dimension of $\SL_{3} (\Z)$ is 3, which
means that the retract $W$ is a 3-dimensional cell complex.  The
closures of the top-dimensional cells in $W$, which are in bijection
with the \Vor \ cells of type $c_{2}$, are homeomorphic to solid cubes
truncated along two pairs of opposite corners (Figure
\ref{soulecube}).  To compute this, one must see how many \Vor cells
of a given type contain a fixed cell of type $c_{2}$ (since the
inclusions of cells in $W$ are the \defn{opposite} of those in $\sC$).

A table of the incidence relations between the cells of $\sC$ and $W$
is given in Table \ref{incidence}.  To interpret the table, let $m =
m(X,Y)$ be the integer in row $X$ and column $Y$.
\begin{itemize}
\item If $m$ is below the diagonal, then the boundary of a cell
of type $Y$ contains $m$ cells of type $X$.
\item If $m$ is above the diagonal, then a cell of type $Y$
appears in the boundary of $m$ cells of type $X$.
\end{itemize}
For instance, the entry $16$ in row $c_{5}$ and column $c_{2}$ means
that a \Vor cell of type $c_{2}$ meets the boundaries of $16$ cells
of type $c_{5}$.  This is the same as the number of vertices in the
Soul\'e cube (Figure
\ref{soulecube}).  Investigation of the table shows that the triangular
(respectively, hexagonal) faces of the Soul\'e cube correspond to the
\Vor \ cells of type $c_{3a}$ (resp., $c_{3b}$).

Figure \ref{souleschlegel} shows a Schlegel diagram for the Soul\'e
cube.  One vertex is at infinity; this is indicated by the arrows on
three of the edges.  This Soul\'e cube is dual to the \Vor cell $C$ of
type $c_{2}$ with minimal vectors given by the columns of the identity
matrix.  The labels on the $2$-faces are additional minimal vectors
that show which \Vor cells contain $C$.  For example, the central
triangle labelled with $(1,1,1)^{t}$ is dual to the \Vor cell of type
$c_{3a}$ with minimal vectors given by those of $C$ together with
$(1,1,1)^{t}$.  Cells of type $c_{4}$ containing $C$ in their closure
correspond to the edges of the figure; the minimal vectors for a given
edge are those of $C$ together with the two vectors on the $2$-faces
containing the edge.  Similarly, one can read off the minimal vectors
of the top-dimensional \Vor cells containing $C$, which correspond to
the vertices of Figure \ref{souleschlegel}.


\begin{table}[htb]
\begin{center}
\begin{tabular*}{0.75\textwidth}{@{\extracolsep{\fill}}c||c|c|c|c|c}
&$c_{5}$&	$c_{4}$&	$c_{3a}$&	$c_{3b}$&	$c_{2}$\\
\hline
\hline
$c_{5}$&	$\bullet$&	$2$&	$3$&	$6$&	$16$\\
$c_{4}$&	$6$&	$\bullet$&	$3$&    $6$&	$24$\\
$c_{3a}$&	$3$&	$1$&	$\bullet$&	$\bullet$&	$4$\\
$c_{3b}$&	$12$&	$4$&	$\bullet$&	$\bullet$&	$6$\\
$c_{2}$&	$12$&	$8$&	$4$&	$3$&	$\bullet$\\
\end{tabular*}
\end{center}
\medskip
\caption{Incidence relations in the \Vor decomposition and the retract
for $\SL_{3} (\Z)$.\label{incidence}}
\end{table}

\begin{figure}[htb]
\begin{center}
\includegraphics[scale=0.35]{diagrams/soulecube}
\end{center}
\caption{\label{soulecube}The Soul\'e cube.}
\end{figure}

\begin{figure}[htb]
\psfrag{111}{$(1,1,1)^t$}
\psfrag{-111}{$(-1,1,1)^t$}
\psfrag{1-11}{$(1,-1,1)^t$}
\psfrag{11-1}{$(1,1,-1)^t$}
\psfrag{110}{$(1,1,0)^t$}
\psfrag{1-10}{$(1,-1,0)^t$}
\psfrag{011}{$(0,1,1)^t$}
\psfrag{01-1}{$(0,1,-1)^t$}
\psfrag{101}{$(1,0,1)^t$}
\psfrag{10-1}{$(1,0,-1)^t$}
\begin{center}
\includegraphics[scale=0.5]{diagrams/souleschlegel}
\end{center}
\caption{A Schlegel diagram of a Soul\'e cube, showing the minimal
vectors that correspond to the $2$-faces.\label{souleschlegel}}
\end{figure}

\subsection{}\label{ss:cohocompnohecke}
%
Now let $p$ be a prime, and let $\Gamma = \Gamma_{0} (p)\subset
\SL_{3} (\Z)$ be the Hecke subgroup of matrices with bottom row
congruent to $(0,0,*) \mod p$ (Example \ref{ex:cong.subgp}).  The
virtual cohomological dimension of $\Gamma$ is $3$, and the cusp
cohomology with constant coefficients can appear in degrees $2$ and
$3$.  One can show that the cusp cohomology in degree $2$ is dual to
that in degree $3$, so for computational purposes it suffices to focus
on degree $3$.  

In terms of $W$, these will be cochains supported on the 3-cells.
Unfortunately we cannot work directly with the quotient $\Gamma
\backslash W$ since $\Gamma$ has torsion: there will be cells taken to
themselves by the $\Gamma$-action, and thus the cells of $W$ need to
be subdivided to induce the structure of a cell complex on $\Gamma
\backslash W$.  Thus when $\Gamma$ has torsion, the ``set of $3$-cells
modulo $\Gamma$'' unfortunately makes no sense.

To circumvent this problem, one can mimic the idea of Manin symbols.
The quotient $\Gamma \backslash \SL_{3} (\Z)$ is in bijection with the
finite projective plane $\Proj^{2} (\F_{p})$, where $\F_p$ is the
field with $p$ elements (cf. Proposition \ref{prop:p1bij}).
The group $\SL_{3} (\Z)$ acts transitively on the set of all $3$-cells
of $W$; if we fix one such cell $w$, its stabilizer $\Stab (w) =
\{\gamma \in \SL_{3} (\Z) \mid \gamma w = w \}$ is a finite subgroup
of $\SL_{3} (\Z)$.  Hence the set of $3$-cells modulo $\Gamma$ should
be interpreted as the set of orbits in $\Proj^{2} (\F_{p})$ of the
finite group $\Stab (w)$.  This suggests describing $H^{3}(\Gamma ;
\C)$ in terms of the space $\sS $ of complex-valued functions $f\colon
\Proj^{2} (\F_{p})\rightarrow \C$.  To carry this out, there are two problems:
\begin{enumerate}
\item How do we explicitly describe $H^{3}(\Gamma ;
\C)$ in terms of $\sS$?
\item How can we isolate the cuspidal subspace $H^{3}_{\text{cusp}}
(\Gamma ;\C) \subset H^{3} (\Gamma ; \C)$ in terms of our description?
\end{enumerate}
Fully describing the solutions to these problems is rather
complicated.  We content ourselves with presenting the following
theorem, which collects together several statements in \cite{agg}.
This result should be compared to Theorems \ref{thm:manin_present} and
\ref{thm:mansym}.

\begin{theorem}[Theorem 3.19 and Summary 3.23 of \cite{agg}]
\label{thm:agg}
  We have 
\[
\dim H^{3} (\Gamma_{0} (p); \C) = \dim H^{3}_{\text{cusp}} (\Gamma_{0}
(p);\C) + 2S_{p},
\]
where $S_{p}$ is the dimension of the space of weight $2$ holomorphic
cusp forms on $\Gamma_{0} (p)\subset \SL_{2} (\Z)$.  Moreover, 
the cuspidal cohomology $H^{3}_{\text{cusp}} (\Gamma_{0}
(p); \C)$ is isomorphic to the vector space of functions $f\colon \Proj^{2}
(\F_p)\rightarrow \C$ satisfying
\begin{enumerate}
\item $f (x,y,z) = f (z,x,y) = f (-x,y,z) = -f (y,x,z)$,
\item $f (x,y,z) + f (-y,x-y,z) + f (y-x,-x,z) = 0$,
\item $f (x,y,0)=0$, and 
\item $\sum_{z=1}^{p-1} f (x,y,z) = 0$.
\end{enumerate}
\end{theorem}

Unlike subgroups of $\SL_{2} (\Z)$, cuspidal cohomology is
apparently much rarer for $\Gamma_{0} (p)\subset \SL_{3} (\Z)$.  The
computations of \cite{agg, fourdutch} show that the only prime levels
$p\leq 337$ with nonvanishing cusp cohomology are $53$, $61$, $79$, $89$,
and $223$.  In all these examples, the cuspidal subspace is
$2$-dimensional. 

For more details of how to implement such computations, we refer to
\cite{agg, fourdutch}.  For further details about the additional
complications arising for higher rank groups, in particular subgroups
of $\SL_{4} (\Z)$, see \cite[Section~3]{computation}.


%%%%%%%%%%%%%%%%%%%%%%%%%%%%%%%
%                             %
%  hecke and modular symbols  %
%                             %
%%%%%%%%%%%%%%%%%%%%%%%%%%%%%%%

\section{Hecke Operators and Modular Symbols}\label{mod.symb.section}
%
\subsection{}
%
There is one ingredient missing so far in our discussion of the
cohomology of arithmetic groups, namely the Hecke operators.\index{Hecke operator}
These are an essential tool in the study of modular forms.  Indeed,
the forms with the most arithmetic significance are the Hecke
eigenforms, and the connection with arithmetic is revealed by the
Hecke eigenvalues. 

In higher rank the situation is similar.  There is an algebra of
Hecke operators acting on the cohomology spaces $H^{*} (\Gamma ;
\MMM)$. The eigenvalues of these operators are conjecturally related
to certain representations of the Galois group.  Just as in the case
$G=\SL_{2} (\R)$, we need tools to compute the Hecke action.

In this section we discuss this problem.  We begin with a general
description of the Hecke operators and how they act on cohomology.
Then we focus on one particular cohomology group, namely the top
degree $H^{\nu} (\Gamma ;\C)$, where $\nu =\vcd (\Gamma)$ and $\Gamma$
has finite index in $\SL_{n} (\Z)$.  This is the setting that
generalizes the modular symbols method from Chapter \ref{ch:modsym}.
We conclude by giving examples of Hecke eigenclasses in the cuspidal
cohomology of $\Gamma_{0} (p)\subset \SL_{3} (\Z)$.


%\subsection{Hecke operators}\label{hecke.ops}
\subsection{}\label{hecke.ops}
%
Let $g\in \SL_{n} (\Q)$.   The group $\Gamma ' = \Gamma \cap
g^{-1}\Gamma g$ has finite index in both $\Gamma$ and $g^{-1}\Gamma
g$.  The element $g$ determines a diagram $C (g)$
\[
\xymatrix{&{\Gamma'\backslash X}\ar[dl]_{s}\ar[dr]^{t}&\\
 {\Gamma\backslash X}&&{\Gamma\backslash X}}
\]
called a \defn{Hecke correspondence}.  The map $s$ is induced by the
inclusion $\Gamma '\subset \Gamma$, while $t$ is induced by the
inclusion $\Gamma '\subset g^{-1}\Gamma g$ followed by the
diffeomorphism $g^{-1}\Gamma g\backslash X\rightarrow \Gamma
\backslash X$ given by left multiplication by $g$.  Specifically, 
\[
s (\Gamma 'x) = \Gamma x, \quad t (\Gamma 'x) = \Gamma gx, \quad x\in X.
\]

The maps $s$ and $t$ are finite-to-one, since the indices $[\Gamma
':\Gamma]$ and $[\Gamma ':g^{-1}\Gamma g]$ are finite.   This implies
that we obtain maps on cohomology 
\[
s^{*}\colon H^{*} (\Gamma \backslash X) \rightarrow
H^{*} (\Gamma' \backslash X), \quad t_{*}\colon H^{*} (\Gamma' \backslash X) \rightarrow
H^{*} (\Gamma \backslash X).
\]
Here the map $s^{*}$ is the usual induced map on cohomology, while the
``wrong-way'' map\footnote{Under the identification $H^{*} (\Gamma
\backslash X; \widetilde{\MMM})\simeq H^{*} (\Gamma ;\MMM)$, the map
$t_{*}$ becomes the transfer map in group cohomology
\cite[III.9]{brown}.} $t_{*}$ is given by summing a class over the finite
fibers of $t$.  These maps can be composed to give a map
\[
T_{g} := t_{*}s^{*}\colon H^{*} (\Gamma \backslash X;
\widetilde{\MMM})\longrightarrow  H^{*} (\Gamma \backslash X;
\widetilde{\MMM}).
\]
This is called the \defn{Hecke operator} associated to $g$.  There is
an obvious notion of isomorphism of Hecke correspondences.  One can
show that up to isomorphism, the correspondence $C (g)$ and thus the
Hecke operator $T_{g}$ depend only on the double coset $\Gamma
g\Gamma$.  One can compose Hecke correspondences, and thus we obtain
an algebra of operators acting on the cohomology, just as in the
classical case.

\begin{example}
Let $n=2$, and let $\Gamma = \SL_{2} (\Z)$.  If we take $g=\diag
(1,p)$, where $p$ is a prime, then the action of $T_{g}$ on $H^{1}
(\Gamma ; M_{k-2})$ is the same as the action of the classical Hecke
operator $T_{p}$ on the weight $k$ holomorphic modular forms.  If we
take $\Gamma = \Gamma_{0} (N)$, we obtain an operator $T (p)$ for all
$p$ prime to $N$, and the algebra of Hecke operators coincides with
the (semisimple) Hecke algebra generated by the $T_{p}$, $(p,N)=1$.
For $p|N$, one can also describe the $U_{p}$ operators in this language.
\end{example}

\begin{example}\label{ex:heckeops}
Now let $n>2$ and let $\Gamma = \SL_{n} (\Z)$. The picture is very
similar, except that now there are several Hecke operators attached to
any prime $p$.  In fact there are $n-1$ operators $T (p,k)$,
$k=1,\dotsc ,n-1$.  The operator $T (p,k)$ is associated to the
correspondence $C (g)$, where $g = \diag (1,\dotsc ,1,p,\dotsc ,p)$
and where $p$ occurs $k$ times.  If we consider the congruence
subgroups $\Gamma_{0} (N)$, we have operators $T (p,k)$ for
$(p,N)=1$ and analogues of the $U_{p}$ operators for $p|N$.

Just as in the classical case, any double coset $\Gamma g \Gamma$ can
be written as a disjoint union of left cosets 
\[
\Gamma g\Gamma  = \coprod_{h\in \Omega} \Gamma h
\]
for a certain finite set of $n\times n$ integral matrices $\Omega$.
For the operator $T (p,k)$, the set $\Omega$ can be taken to be all
upper-triangular matrices of the form \cite[Proposition 7.2]{krieg}
\[
\left(\begin{array}{cccc}
p^{e_{1}}&&a_{ij}\\
&\ddots&&\\
&&p^{e_{n}}&\\
\end{array} \right),
\]
where
\begin{itemize}
\item $e_{i}\in \{0,1 \}$ and exactly $k$ of the $e_{i}$ are equal to
$1$ and
\item $a_{ij}=0$ unless $e_{i}=0$ and $e_{j}=1$, in which case
$a_{ij}$ satisfies $0\leq a_{ij}<p$.
\end{itemize}


\end{example}

\begin{remark}
The number of coset representatives for the operator $T (p,k)$ is the
same as the number of points in the finite Grassmannian $G (k,n)
(\F_{p})$.  A similar phenomenon is true for the Hecke operators for
any group $G$, although there are some subtleties \cite{gross}.  
\end{remark}



%\subsection{Modular symbols}\label{mod.symbs}
\subsection{}\label{mod.symbs}
%
Recall that in Section~\ref{wrr.section} we constructed the \Vor
decomposition $\sC$ and the well-rounded retract $W$ and that we can
use them to compute the cohomology $H^{*} (\Gamma ;\MMM)$.
Unfortunately, we cannot directly use them to compute the action of the
Hecke operators on cohomology, since the Hecke operators do not act
cellularly on $\sC$ or $W$.  The problem is that the Hecke image of a
cell in $\sC$ (or $W$) is usually not a union of cells in $\sC$ (or
$W$).  This is already apparent for $n=2$.  The edges of $\sC$ are the
$\SL_{2} (\Z)$-translates of the ideal geodesic $\tau$ from $0$ to
$\infty$ (Example \ref{sl2}).  Applying a Hecke operator takes such an
edge to a union of ideal geodesics, each with vertices at a pair of
cusps.  In general such geodesics are not an $\SL_{2} (\Z)$-translate
of~$\tau$.  

For $n=2$, one solution is to work with all possible ideal geodesics
with vertices at the cusps, in other words the space of modular
symbols $\M_{2}$ from Section~\ref{sec:modsym}.  Manin's trick (Proposition
\ref{prop:manintrick}) shows how to write any modular symbol as a
linear combination of unimodular symbols, by which we mean modular
symbols supported on the edges of $\sC$.  These are the ideas we now
generalize to all $n$.

\begin{definition}\label{def:modular.symbols}
Let $S_{0}$ be the $\Q$-vector space
spanned by the symbols $\vv = [v_{1},\dotsc ,v_{n}]$, where $v_{i}\in
\Q^{n}\smallsetminus \{0 \}$, modulo the following relations:
\begin{enumerate}
\item If $\tau $ is a permutation on $n$ letters, then
\[
[v_{1},\ldots,v_{n}] = \sign (\tau ) [\tau (v_{1}),\ldots,\tau
(v_{n})], 
\]
where $\sign (\tau )$ is the sign of $\tau $.
\item  \label{projpts} If $q\in \Q^{\times }$, then 
\[
[q v_{1},v_{2}\ldots,v_{n}] = [v_{1},\ldots,v_{n}].
\]
\item \label{lindep} If the points $v_{1},\dotsc ,v_{n}$ are linearly dependent, then $\vv = 0$.
\end{enumerate}
Let $B\subset S_{0}$ be the subspace generated by linear combinations
of the form 
\begin{equation}\label{fundrel}
\sum _{i=0}^{n}  (-1)^{i}
[v_{0},\ldots,\hat{v_{i}},\ldots,v_{n}],  
\end{equation}
where $v_{0},\dotsc ,v_{n}\in
\Q^{n}\smallsetminus \{0 \}$ and where $\hat{v_{i}}$ means to omit $v_{i}$.
\end{definition}

We call $S_{0}$ the space of \defn{modular symbols}.  We caution the
reader that there are some differences in what we call modular symbols
and those found in Section~\ref{sec:modsym} and Definition
\ref{defn:modsym}; we compare them in 
Section~\ref{sec:comparison}.  The group $\SL_{n} (\Q)$ acts on $S_{0}$ by
left multiplication: $g\cdot \vv = [gv_{1},\dotsc ,gv_{n}]$.  This
action preserves the subspace $B$ and thus induces an action on the
quotient $M = S_{0}/B$.  For $\Gamma \subset \SL_{n} (\Z)$ a finite
index subgroup, let $M_{\Gamma}$ be the space of $\Gamma$-coinvariants
in $M$.  In other words, $M_{\Gamma}$ is the quotient of $M$ by the
subspace generated by $\{ m - \gamma \cdot m \mid \gamma \in
\Gamma\}$.

The relationship between modular symbols and the cohomology of
$\Gamma$ is given by the following theorem, first proved for $\SL_{n}$
by Ash and Rudolph \cite{ash.rudolph} and by Ash for general $G$
\cite{ash.minmod}:

\begin{theorem}[\cite{ash.minmod, ash.rudolph}]\label{thm:modsymbiso}
 Let $\Gamma \subset \SL_{n} (\Z)$ be a
finite index subgroup.  There is an isomorphism
\begin{equation}\label{modsymbiso}
M_{\Gamma}\xrightarrow{\sim} H^{\nu} (\Gamma ; \Q),
%\xymatrix@1{S_{0,\Gamma}\ar[r]^{\sim}& H^{\nu} (\Gamma ; \Q)}
\end{equation}
where $\Gamma$ acts trivially on $\Q$ and where $\nu = \vcd (\Gamma )$.
\end{theorem}

We remark that Theorem \ref{thm:modsymbiso} remains true if $\Q$ is
replaced with nontrivial coefficients as in Section \ref{ss:coho}.
Moreover, if $\Gamma$ is assumed to be torsion-free then we can
replace $\Q$ with $\Z$.


The great virtue of $M_{\Gamma}$ is that it admits an action of the
Hecke operators.  Given a Hecke operator $T_{g}$, write 
the double coset $\Gamma g\Gamma$ as a disjoint union of left cosets
\begin{equation}\label{eq:coset.decomp}
\Gamma g\Gamma  = \coprod_{h\in \Omega} \Gamma h
\end{equation}
as in Example \ref{ex:heckeops}.  
Any class in $M_{\Gamma}$ can be lifted to a representative $\eta =
\sum q (\vv)\vv \in S_0$, where $q (\vv)\in \Q$
and almost all $q (\vv)$ vanish.  
Then we define
\begin{equation}\label{eq:heckeaction}
T_{g} (\vv) = \sum_{h\in \Omega} h\cdot \vv
\end{equation}
and extend to $\eta$ by linearity.  The right side of
\eqref{eq:heckeaction} depends on the choices of $\eta$ and $\Omega$,
but after taking quotients and coinvariants, we obtain a well-defined
action on cohomology via \eqref{modsymbiso}.

\subsection{}\label{sec:comparison}
%
The space $S_{0}$ is closely related to the space $\M_{2}$ from
Section~\ref{sec:modsym} and Section~\ref{sec:modsymgen}.  Indeed, $\M_{2}$ was
defined to be the quotient $(F/R)/ (F/R)_{\text{tor}}$, where $F$ is
the free abelian group generated by ordered pairs
\begin{equation}\label{symbol.def}
\{\alpha ,\beta \}, \quad \alpha ,\beta \in \Proj^{1} (\Q),
\end{equation}
and $R$ is the subgroup generated by elements of the form
\begin{equation}\label{fund.rel}
  \{\alpha,\beta\} + \{\beta,\gamma\} + \{\gamma,\alpha\}, \quad
\alpha ,\beta, \gamma  \in \Proj^{1} (\Q).
\end{equation}
The only new feature in Definition \ref{def:modular.symbols} is item
\eqref{lindep}.  For $n=2$ this corresponds to the condition $\{\alpha,\alpha
\}=0$, which follows from \eqref{fund.rel}.  
We have
\[
S_{0}/B \simeq \M_{2} \otimes \Q.
\]
Hence there are two differences between $S_{0}$ and $\M_{2}$: our
notion of modular symbols uses rational coefficients instead of
integral coefficients and is the space of symbols \emph{before}
dividing out by the subspace of relations $B$; we further caution the
reader that this is somewhat at
odds with the literature. 

We also remark that the general arbitrary weight definition of modular
symbols for a subgroup $\Gamma \subset \SL_{2} (\Z)$ given in
Section~\ref{sec:modsymgen} also includes taking $\Gamma$-coinvariants, as
well as extra data for a coefficient system.  We have not included the
latter data since our emphasis is trivial coefficients, although it
would be easy to do so in the spirit of Section~\ref{sec:modsymgen}.

Elements of $\M_{2}$ also have a geometric interpretation: the symbol
$\{\alpha ,\beta \}$ corresponds to the ideal geodesic in $\fh$ with
endpoints at the cusps $\alpha$ and $\beta$.  We have a similar
picture for the symbols $\vv = [v_{1},\dotsc ,v_{n}]$.  We can assume
that each $v_{i}$ is primitive, which means that each $v_{i}$
determines a vertex of the \Vor polyhedron $\Pi$.  The rational cone
generated by these vertices determines a subset $\Delta (\vv)\subset
D$, where $D$ is the linear model of the symmetric space $X = \SL_{n}
(\R)/\SO (n)$ from Section~\ref{vcd.section}.  This subset $\Delta (\vv)$ is
then an ``ideal simplex'' in $X$.  There is also a connection between
$\Delta (\vv)$ and torus orbits in $X$; we refer to \cite{ash.minmod}
for a related discussion.

%\subsection{The modular symbol  algorithm}
\subsection{}
%
Now we need a generalization of the Manin trick
(Section~\ref{sec:manintrick}).  This is known in the literature as the
\defn{modular symbols algorithm}.

We can define a kind of norm function on $S_{0}$ as follows.  Let $\vv
=[v_{1},\dotsc ,v_{n}]$ be a modular symbol.  For each $v_{i}$, 
choose $\lambda_{i}\in \Q^{\times}$ such that $\lambda_{i}v_{i}$ is
primitive.  Then we define 
\[
\norm{\vv} := |\det (\lambda_{1}v_{1},\dotsc ,\lambda_{n}v_{n})| \in \Z .
\]
Note that $\norm{\vv}$ is well defined, since the $\lambda_{i}$ are
unique up to sign, and permuting the $v_{i}$ only changes the
determinant by a sign.  We extend $\norm{\phantom{a}}$ to all of $S_{0}$
by taking the maximum of $\norm{\phantom{a}}$ over the support of any
$\eta \in S_{0}$: if $\eta =\sum q (\vv)\vv$, where $q (\vv)\in \Q$
and almost all $q (\vv)$ vanish, then we put
\[
\norm{\eta} = \Max_{q (\vv) \not =0} \norm{\vv}.
\]
We say a modular symbol $\eta $ is \defn{unimodular} if $\norm{\eta }
= 1$.  It is clear that the images of the unimodular symbols generate
a finite-dimensional subspace of $M_{\Gamma}$.  The next theorem shows
that this subspace is actually \emph{all} of $M_{\Gamma }$.

\begin{theorem}[\cite{ash.rudolph, barvinok}]
\label{thm:modsymbalg}
The space $M_{\Gamma }$ is spanned by
the images of the unimodular symbols.  More precisely, given any symbol $\vv\in S_{0}$ with
$\norm{\vv}>1$, 
\begin{enumerate}
\item in $S_{0}/B$ we may write 
\begin{equation}\label{relation}
\vv = \sum q (\ww)\ww , \quad q (\ww)\in \Z ,
\end{equation}
where if $q (\ww)\not =0$, then $\norm{\ww} = 1$, and 
\item the number of terms on the right side of \eqref{relation} is bounded
by a polynomial in $\log \|\vv\|$ that depends only on the dimension $n$.
\end{enumerate}
\end{theorem}

\begin{proof}
%
(Sketch)
Given a modular symbol $\vv =[v_{1},\dotsc ,v_{n}]$, we may assume
that the points $v_{i}$ are primitive.  We will show that if
$\norm{\vv}>1$, we can find a point $u$ such that when we apply the
relation \eqref{fundrel} using the points $u,v_{1},\dotsc ,v_{n}$, all
terms other than $\vv$ have norm less than $\norm{\vv}$.  We call such
a point a \defn{reducing point} for $\vv$.

Let $P\subset \R^{n}$ be the open parallelotope
\[
P:=\Bigl\{\sum \lambda _{i}v _{i}\Bigm | |\lambda_{i} |< \norm{\vv }^{-1/n}
\Bigr\}. 
\]
Then $P$ is an $n$-dimensional centrally symmetric convex body with
volume $2^{n}$.  By Minkowski's theorem from the geometry of numbers
(cf.~\cite[IV.2.6]{fr.tay}), $P\cap \Z ^{n}$ contains a nonzero point
$u$.  
Using \eqref{fundrel}, we find 
\begin{equation}\label{move}
\vv = \sum_{i=1}^{n} (-1)^{i-1}\vv_{i} (u),
\end{equation}
where $\vv_{i} (u)$ is the symbol
\[
\vv_{i} (u) = [v_{1},\dotsc ,v_{i-1},u, v_{i+1},\dotsc ,v_{n}].
\]
Moreover, it is easy to see that the new symbols satisfy
\begin{equation}\label{barvest}
0\leq \norm{\vv_{i} (u)}<\norm{\vv}^{(n-1)/n}, \quad i=1,\dotsc ,n.
\end{equation}
This completes the proof of the first statement.

To prove the second statement, we must estimate how many times
relations of the form \eqref{move} need to be applied to obtain
\eqref{relation}.  A nonunimodular symbol produces at most $n$ new
modular symbols after \eqref{move} is performed; we potentially have
to apply \eqref{move} again to each of the symbols that result, which
in turn could produce as many as $n$ new symbols for each.  Hence we
can visualize the process of constructing \eqref{relation} as building
a rooted tree, where the root is $\vv$, the leaves are the symbols
$\ww$, and where each node has at most $n$ children.  It is not hard
to see that the bound \eqref{barvest} implies that the depth of this
tree (i.e., the longest length of a path from the root to a leaf) is
$O(\log \log \norm{\vv})$.  From this the second statement follows
easily.
%
\end{proof}

Statement (1) of Theorem \ref{thm:modsymbalg} is due to Ash and Rudolph
\cite{ash.rudolph}.  Instead of $P$, they used the larger
parallelotope $P'$ defined by 
\[
P':=\Bigl\{\sum \lambda _{i}v _{i}\Bigm | |\lambda_{i} |< 1
\Bigr\},
\]
which has volume $2^{n}\norm{\vv}$.  The observation that $P'$ can be
replaced by $P$ and the proof of (2) are both due to Barvinok
\cite{barvinok}.

\subsection{}\label{sec:retractandmanin}
%
The relationship between Theorem \ref{thm:modsymbalg} and Manin's
trick should be clear.  For $\Gamma \subset \SL_{2} (\Z)$, the Manin
symbols correspond exactly to the unimodular symbols mod $\Gamma$.  So
Theorem \ref{thm:modsymbalg} implies that every modular symbol (in the
language of Section~\ref{sec:modsymgen}) is a linear combination of Manin
symbols.  This is exactly the conclusion of Proposition
\ref{prop:mangen}.

In higher rank the relationship between Manin symbols and unimodular
symbols is more subtle.  In fact there are two possible notions of
``Manin symbol,'' which agree for $\SL_{2} (\Z)$ but not in general.
One possibility is the obvious one: a Manin symbol is a unimodular
symbol.  

The other possibility is to define a Manin symbol to be a modular
symbol corresponding to a top-dimensional cell of the retract $W$.
But for $n\geq 5$, such modular symbols need not be unimodular.  In
particular, for $n=5$ there are two equivalence classes of
top-dimensional cells.  One class corresponds to the unimodular
symbols, the other to a set of modular symbols of norm $2$.  However,
Theorems~\ref{thm:modsymbiso} and~\ref{thm:modsymbalg} show that
$H^{\nu} (\Gamma ; \Q)$ is spanned by unimodular symbols.  Thus as far
as this cohomology group is concerned, the second class of symbols is
in some sense unnecessary.


\subsection{}\label{ss:cohocomphecke}
%
We return to the setting of Section~\ref{ss:cohocompnohecke} and give
examples of Hecke eigenclasses in the cusp cohomology of $\Gamma =
\Gamma_{0} (p)\subset \SL_{3} (\Z)$.  We closely follow \cite{agg,
fourdutch}.  Note that since the top of the cuspidal range for
$\SL_{3}$ is the same
as the virtual cohomological dimension $\nu$, we can use modular
symbols to compute the Hecke action on cuspidal classes.

Given a prime $l$ coprime to $p$, there are two Hecke operators of
interest $T (l,1)$ and $T (l,2)$.  We can compute the action of these
operators on $H^{3}_{\text{cusp}} (\Gamma ; \C)$ as follows.  Recall
that $H^{3}_{\text{cusp}} (\Gamma ; \C)$ can be identified with a
certain space of functions $f\colon \Proj^{2} (\F_{p})\rightarrow \C$
(Theorem \ref{thm:agg}).  Given $x\in \Proj^{2} (\F_p)$, let $Q_{x}\in
\SL_{3} (\Z)$ be a matrix such that $Q_{x}\mapsto x$ under the
identification $\Proj^{2} (\F_p) \xrightarrow{\sim} \Gamma \backslash
\SL_{3} (\Z)$.  Then $Q_{x}$ determines a unimodular symbol $[Q_{x}]$
by taking the $v_{i}$ to be the columns of $Q_{x}$.  Given any Hecke
operator $T_{g}$, we can find coset representatives $h_{i}$ such that
$\Gamma g \Gamma  = \coprod \Gamma h_{i}$ (explicit representatives
for $\Gamma = \Gamma_{0} (p)$ and $T_{g} = T (l,k)$ are given in
\cite{agg, fourdutch}).  The modular symbols $[h_{i}Q_{x}]$ are no
longer unimodular in general, but we can apply Theorem
\ref{thm:modsymbalg} to write 
\[
[h_{i}Q_{x}] = \sum_{j} [R_{ij}], \quad R_{ij}\in \SL_{3} (\Z).
\]
Then for $f\colon \Proj^{2} (\F_p)\rightarrow \C$ as in Theorem
\ref{thm:agg}, we have 
\[
(T_{g}f) (x) = \sum_{i,j} f (\overline{R_{ij}}), 
\]
where $\overline{R_{ij}}$ is the class of $R_{ij}$ in $\Proj^{2}
(\F_p)$. 

Now let $\xi\in H^{3}_{\text{cusp}} (\Gamma ; \C)$ be a simultaneous
eigenclass for all the Hecke operators $T (l,1)$, $T (l,2)$, as $l$
ranges over all primes coprime with $p$.  General considerations from
the theory of automorphic forms imply that the eigenvalues $a (l,1)$,
$a (l,2)$ are complex conjugates of one other.  Hence it suffices to
compute $a (l,1)$.  We give two examples
of cuspidal eigenclasses for two different prime levels.

\begin{example}
Let $p=53$.  Then $H^{3}_{\text{cusp}} (\Gamma_{0} (53); \C)$ is
$2$-dimensional.  Let $\eta = (1+\sqrt{-11})/2$.  One eigenclass is
given by the data 
\begin{center}
\begin{tabular*}{0.75\textwidth}{@{\extracolsep{\fill}}c||c|c|c|c|c|c}
$l$&2&3&5&7&11&13\\
\hline
$a (l,1)$&$-1-2\eta$&$-2+2\eta$&$1$&$-3$&$1$&$-2-12\eta$\\
\end{tabular*}
\end{center}
and the other is obtained by complex conjugation.
\end{example}

\begin{example}
Let $p=61$.  Then $H^{3}_{\text{cusp}} (\Gamma_{0} (61); \C)$ is
$2$-dimensional.  Let $\omega  = (1+\sqrt{-3})/2$.  One eigenclass is
given by the data 
\begin{center}
\begin{tabular*}{0.9\textwidth}{@{\extracolsep{\fill}}c||c|c|c|c|c|c}
$l$&2&3&5&7&11&13\\
\hline
$a (l,1)$&$1-2\omega$&$-5+4\omega $&$-2+4\omega $&$-6\omega $&$-2+2\omega $&$-2-4\omega$\\
\end{tabular*}
\end{center}
and the other is obtained by complex conjugation.
\end{example}

%%%%%%%%%%%%%%%%%%%%%%%%%%%%%
%                           %
%  other cohomology groups  %
%                           %
%%%%%%%%%%%%%%%%%%%%%%%%%%%%%

\section{Other Cohomology Groups}\label{other.coho}
%
\subsection{}
%
In Section~\ref{mod.symb.section} we saw how to compute the Hecke action on
the top cohomology group $H^{\nu} (\Gamma ; \C)$.  Unfortunately for
$n\geq 4$, this cohomology group does not contain any cuspidal
cohomology.  The first case is $\Gamma \subset \SL_{4} (\Z )$; we have
$\vcd (\Gamma) = 6$, and the cusp cohomology lives in degrees $4$ and
$5$.  One can show that the cusp cohomology in degree $4$ is dual to
that in degree $5$, so for computational purposes it suffices to be
able to compute the Hecke action on $H^{5} (\Gamma ; \C)$.  But
modular symbols do not help us here.

In this section we describe a technique to compute the Hecke action on
$H^{\nu -1} (\Gamma ; \C)$, following \cite{experimental}.  The
technique is an extension of the modular symbol algorithm to these
cohomology groups.  In principle the ideas in this section can be
modified to compute the Hecke action on other cohomology groups
$H^{\nu -k} (\Gamma ; \C)$, $k>1$, although this has not been
investigated\footnote{The first interesting case is $n=5$, for which
the cuspidal cohomology lives in $H^{\nu -2}$.}.  For $n=4$, we have
applied the algorithm in joint work with Ash and McConnell to
investigate computationally the cohomology $H^{5} (\Gamma ; \C)$,
where $\Gamma_{0} (N)\subset \SL_{4} (\Z)$ \cite{computation}.

%\subsection{The Sharbly complex}\label{sharbly.sect}
\subsection{}\label{sharbly.sect}
%
To begin, we need an analogue of Theorem \ref{thm:modsymbiso} for
lower degree cohomology groups.  In other words, we need a
generalization of the modular symbols for other cohomology groups.
This is achieved by the \defn{sharbly complex} $S_{*}$:

\begin{definition}[\cite{ash.sharb}] \label{sharbly.complex}
 Let 
$\left\{S_{*},\partial \right\}$ be the chain complex given by the following data:
\begin{enumerate}
\item For $k\geq 0$, $S_{k}$ is the $\Q$-vector space generated
by the symbols $\uu  = [v_{1},\ldots,v_{n+k}]$, where $v_{i}\in
\Q^{n}\smallsetminus \{0 \}$, modulo the relations:
\begin{enumerate}
\item If $\tau $ is a permutation on $(n+k)$ letters, then
\[
[v_{1},\ldots,v_{n+k}] = \sign (\tau ) [\tau (v_{1}),\ldots,\tau
(v_{n+k})], 
\]
where $\sign (\tau )$ is the sign of $\tau $.
\item  If $q \in \Q^{\times }$, then 
\[
[q v_{1},v_{2}\ldots,v_{n+k}] = [v_{1},\ldots,v_{n+k}].
\]
\item If the rank of the matrix
$(v_{1},\ldots,v_{n+k})$ is less than $n$, then $\uu = 0$.
\end{enumerate}
\item For $k>0$, the boundary map $\partial \colon S_{k}\rightarrow S_{k-1}$ is 
$$[v_{1},\ldots,v_{n+k}] \longmapsto \sum _{i=1}^{n+k}  (-1)^{i}
[v_{1},\ldots,\hat{v_{i}},\ldots,v_{n+k}].  $$
We define $\partial$ to be identically zero on $S_{0}$.
\end{enumerate}
\end{definition}

The elements $$\uu = [v_{1},\dots ,v_{n+k}]$$ are called
\defn{$k$-sharblies}\footnote{The terminology for $S_{*}$ is due to
Lee Rudolph, in honor of Lee and Szczarba.  They introduced a very
similar complex in \cite{l.s} for $\SL_{3} (\Z)$.}.  The $0$-sharblies
are exactly the modular symbols from Definition
\ref{def:modular.symbols}, and the subspace $B\subset S_{0}$ is the
image of the boundary map $\partial \colon S_{1}\rightarrow S_{0}$.

There is an obvious left action of $\Gamma $ on $S_{* }$ commuting
with $\partial $.  For any $k\geq 0$, let $S_{k,\Gamma }$ be the space
of $\Gamma $-coinvariants.  Since the boundary map $\partial$ commutes
with the $\Gamma $-action, we obtain a complex $(S_{*,\Gamma},
\partial_{\Gamma})$.  The following theorem shows that this complex
computes the cohomology of $\Gamma $:

\begin{theorem}[\cite{ash.sharb}]
\label{thm:sharbiso}
There is a natural isomorphism 
\[
H^{\nu - k}(\Gamma ; \C) \xrightarrow{\sim} H_{k}
(S_{*,\Gamma }\otimes \C ).
\]
\end{theorem}

\subsection{}
%
We can extend our norm function $\|\phantom{a}\|$ from modular
symbols to all of $S_{k}$ as follows.  Let $\uu = [v_{1},\dots
,v_{n+k}]$ be a $k$-sharbly, and let $Z(\uu)$ be the set of all
submodular symbols determined by $\uu$.  In other words, $Z (\uu)$
consists of the modular symbols of the form $[v_{i_{1}},\dotsc
,v_{i_{n}}]$, where $\{i_{1},\dotsc ,i_{n} \}$ ranges over all
$n$-fold subsets of $\{1,\dotsc ,n+k \}$.  Define $\|\uu \|$ by
\[
\|\uu\| = \Max_{\vv \in Z (\uu )} \| \vv \|.
\]
Note that $\| \uu \|$ is well defined modulo the relations in
Definition \ref{sharbly.complex}.  As for modular symbols, we extend
the norm to sharbly chains $\xi = \sum q (\uu)\uu $ taking the maximum
norm over the support.  Formally, we let $\support (\xi) = \{\uu \mid
q (\uu)\not =0 \}$ and $Z (\xi) = \bigcup_{\uu \in \support (\xi )} Z
(\uu)$, and then we define $\|\xi \|$ by
\[
\|\xi \| = \Max_{\vv \in Z (\xi )} \| \vv \|.
\]

We say that $\xi $ is \defn{reduced} if $\|\xi \| = 1$.  Hence $\xi$
is reduced if and only if all its submodular symbols are unimodular or
have determinant $0$.  Clearly there are only finitely many reduced
$k$-sharblies modulo $\Gamma$ for any $k$.

In general the cohomology groups $H^{*} (\Gamma ; \C)$ are \emph{not}
spanned by reduced sharblies.  However, it is known (cf. \cite{M91})
that for $\Gamma \subset \SL _{4} (\Z )$, the group $H^{5} (\Gamma ;
\C)$ is spanned by reduced $1$-sharbly cycles.  The best one can say
in general is that for each pair $n,k$, there is an integer $N=N
(n,k)$ such that for $\Gamma \subset \SL_{n} (\Z)$, $H^{\nu -k}
(\Gamma; \C)$ is spanned by $k$-sharblies of norm $\leq N$.  This set
of sharblies is also finite modulo $\Gamma$, although it is not known
how large $N$ must be for any given pair $n,k$.

%\subsection{Sharblies and Voronoi}\label{that.section}
\subsection{}\label{that.section}
%
Recall that the cells of the well-rounded retract $W$ are indexed by
sets of primitive vectors in $\Z^{n}$.  Since any primitive vector
determines a point in $\Q^{n}\smallsetminus \{0 \}$ and since sets of
such points index sharblies, it is clear that there is a close
relationship between $S_{*}$ and the chain complex associated to $W$,
although of course $S_{*}$ is much bigger.  In any case, both
complexes compute $H^{*} (\Gamma ; \C)$.

The main benefit of using the sharbly complex to compute cohomology is
that it admits a Hecke action.  Suppose $\xi = \sum q (\uu )\uu $ is a
sharbly cycle mod $\Gamma $, and consider a Hecke operator $T_{g}$.
Then we have
\begin{equation}\label{how.hecke.acts}
T_{g}(\xi ) = \sum _{h\in \Omega , \uu } n (\uu )h\cdot \uu, 
\end{equation}
where $\Omega$ is a set of coset representatives as in
\eqref{eq:coset.decomp}.  
Since $\Omega \not \subset \SL _{n} (\Z )$ in general, the 
Hecke image of a
reduced sharbly is not usually reduced.

\subsection{}
%
We are now ready to describe our algorithm for the computation of the
Hecke operators on $H^{\nu -1} (\Gamma ;\C )$.   It suffices to describe an
algorithm that takes as input a $1$-sharbly cycle $\xi $ and produces
as output a cycle $\xi '$ with
\begin{enumerate}
\item [(a)] the classes of $\xi $ and $\xi '$ in $H^{\nu -1} (\Gamma
;\C )$ the same, and 
\item [(b)] $\|\xi' \| < \|\xi \|$ if $\|\xi \|>1$.
\end{enumerate}

Below, we will present an algorithm satisfying (a).  In
\cite{experimental}, we conjectured (and presented evidence) that the
algorithm satisfies (b) for $n\leq 4$.  Further evidence is provided
by the computations in \cite{computation}, which relied on the
algorithm to compute the Hecke action on $H^{5} (\Gamma ; \C)$, where
$\Gamma = \Gamma_{0} (N)\subset \SL_{4} (\Z)$.

The idea behind the algorithm is simple: given a $1$-sharbly cycle
$\xi$ that is not reduced, (i) simultaneously apply the modular symbol
algorithm (Theorem \ref{thm:modsymbalg}) to each of its submodular
symbols, and then (ii) package the resulting data into a new $1$-sharbly
cycle.  Our experience in presenting this algorithm is that most
people find the geometry involved in (ii) daunting.  Hence we will
give details only for $n=2$ and will provide a sketch for $n>2$.
Full details are contained in \cite{experimental}.  Note that $n=2$ is
topologically and arithmetically uninteresting, since we are computing
the Hecke action on $H^{0} (\Gamma ; \C)$; nevertheless, the geometry
faithfully represents the situation for all $n$.

\subsection{}\label{sl2.start}
%
Fix $n=2$, let $\xi \in S_{1}$ be a $1$-sharbly cycle mod $\Gamma$ for
some $\Gamma \subset \SL_{2} (\Z )$, and suppose $\xi $ is not
reduced.  Assume $\Gamma$ is torsion-free to simplify the
presentation.  

Suppose first that all submodular symbols $\vv \in Z (\xi )$ are
nonunimodular.  Select reducing points for each $\vv \in Z (\xi )$
and make these choices $\Gamma $-equivariantly.  This means the
following.  Suppose $\uu,\uu '\in \support \xi $ and $\vv \in \support
(\partial \uu )$ and $\vv' \in \support (\partial \uu' )$ are modular
symbols such that $\vv = \gamma\cdot \vv ' $ for some $\gamma \in
\Gamma $.  Then we select reducing points $w$ for $\vv $ and $w'$ for
$\vv '$ such that $w = \gamma\cdot w'$.  (Note that since $\Gamma$ is
torsion-free, no modular symbol can be identified to itself by an
element of $\Gamma$; hence $\vv \not = \vv '$.)  This is possible
since if $\vv $ is a modular symbol and $w$ is a reducing point for
$\vv $, then $\gamma \cdot w $ is a reducing point for $\gamma \cdot
\vv$ for any $\gamma \in \Gamma $.  Because there are only finitely
many $\Gamma $-orbits in $Z (\xi )$, we can choose reducing points
$\Gamma $-equivariantly by selecting them for some set of orbit
representatives.

It is important to note that $\Gamma $-equivariance is the only
global criterion we use when selecting reducing.  In
particular, there is a priori no relationship among the three reducing
points chosen for any $\uu \in \support \xi $.

\subsection{}\label{ss:buildingxiprime}
%
Now we want to use the reducing points and the $1$-sharblies in $\xi$ to
build $\xi '$.  Choose $\uu = [v_{1},v_{2},v_{3}]\in
\support \xi $, and denote the reducing point for $[v_{i},v_{j}]$ by $w_{k}$,
where $\{i,j,k \} = \{1,2,3 \}$.  We use the $v_{i}$ and the
$w_{i}$ to build a $2$-sharbly chain $\eta (\uu )$ as follows.

Let $P$ be an octahedron in $\R ^{3}$.  Label the vertices of $P$ with
the $v_{i}$ and $w_{i}$ such that the vertex labeled $v_{i}$ is
opposite the vertex labeled $w_{i}$ (Figure \ref{octa.fig}).
Subdivide $P$ into four tetrahedra by connecting two opposite
vertices, say $v_{1}$ and $w_{1}$, with an edge (Figure
\ref{octa-hom.fig}).
For each tetrahedron $T$, take the labels of four vertices and arrange
them into a quadruple.  If we orient $P$, then we can use the induced
orientation on $T$ to order the four primitive points.  In this way,
each $T$ determines a $2$-sharbly, and $\eta (\uu)$ is defined to be
the sum.  For example, if we use the decomposition in Figure
\ref{octa-hom.fig}, we have
\begin{equation}\label{eta}
\eta (\uu ) = [v_{1},v_{3},
v_{2},w_{1}]+[v_{1},w_{2},v_{3},w_{1}]+[v_{1},w_{3},w_{2},w_{1}]+[v_{1},v_{2},w_{3},w_{1}].
\end{equation}
Repeat this construction for all $\uu \in \support \xi $, and let 
$\eta =\sum q (\uu ) \eta (\uu )$.  Finally, let $\xi ' = \xi +\partial \eta $.

\begin{figure}[htb]
\begin{center}
\psfrag{v1}{$v_{1}$}
\psfrag{v2}{$v_{2}$}
\psfrag{v3}{$v_{3}$}
\psfrag{w1}{$w_{1}$}
\psfrag{w2}{$w_{2}$}
\psfrag{w3}{$w_{3}$}
\includegraphics[scale=.5]{diagrams/octa}
\end{center}
\caption{\label{octa.fig}}
\end{figure}

\begin{figure}[htb]
\begin{center}
\psfrag{v1}{$v_{1}$}
\psfrag{v2}{$v_{2}$}
\psfrag{v3}{$v_{3}$}
\psfrag{w1}{$w_{1}$}
\psfrag{w2}{$w_{2}$}
\psfrag{w3}{$w_{3}$}
\includegraphics[scale=.5]{diagrams/octa-hom}
\end{center}
\caption{\label{octa-hom.fig}}
\end{figure}

\subsection{}\label{endofdiscussion}
%
By construction, $\xi' $ is a cycle mod $\Gamma$ in the same
class as $\xi $.  We claim in addition that no submodular symbol from
$\xi $ appears in $\xi '$.  To see this, consider $\partial \eta (\uu
)$.  From \eqref{eta}, we have
\begin{multline}\label{bound}
\partial \eta (\uu ) = - [v_{1},v_{2},v_{3}] + [v_{1},v_{2},w_{3}] +
[v_{1},w_{2},v_{3}] + [w_{1},v_{2},v_{3}] \\
- [v_{1},w_{2},w_{3}] -
[w_{1},v_{2},w_{3}] - [w_{1},w_{2},v_{3}] + [w_{1},w_{2},w_{3}].
\end{multline}
Note that this is the boundary in $S_{*}$, not in $S_{*,\Gamma }$.
Furthermore, $\partial \eta (\uu )$ is independent of which pair of
opposite vertices of $P$ we connected to build $\eta (\uu )$.

From \eqref{bound}, we see that in $\xi +\partial \eta $ the
$1$-sharbly $-[v_{1},v_{2},v_{3}]$ is canceled by $\uu \in \support
\xi $.  We also claim that $1$-sharblies in \eqref{bound} of the form
$[v_{i},v_{j},w_{k}]$ vanish in 
$\partial_{\Gamma } \eta $.

To see this, let $\uu $,$\uu '\in \support \xi $, and suppose $\vv =
[v_{1},v_{2}]\in \support \partial \uu $ equals $\gamma \cdot \vv '$
for some $\vv ' = [v_{1}',v_{2}']\in \support \partial \uu '$.  Since
the reducing points were chosen $\Gamma $-equivariantly, we have
$w=\gamma \cdot w'$.  This means that the $1$-sharbly $[v_{1},v_{2},
w]\in \partial \eta (\uu )$ will be canceled mod $\Gamma $ by
$[v_{1}',v_{2}', w']\in \partial \eta (\uu' )$.  Hence, in passing
from $\xi $ to $\xi '$, the effect in $\coinv{*}$ is to replace $\uu $
with \emph{four} $1$-sharblies in $\support \xi '$:
\begin{equation}\label{tfm}
[v_{1},v_{2},v_{3}]\longmapsto - [v_{1},w_{2},w_{3}] -
[w_{1},v_{2},w_{3}] - [w_{1},w_{2},v_{3}] + [w_{1},w_{2},w_{3}].
\end{equation}
Note that in \eqref{tfm}, there are no $1$-sharblies of the form
$[v_{i},v_{j},w_{k}]$.

\begin{remark}\label{imp1}
For implementation purposes, it is not necessary to explicitly
construct $\eta $.  Rather, one may work directly with \eqref{tfm}.
\end{remark}

\subsection{}\label{interior.ms}
%
Why do we expect $\xi ' $ to satisfy $\|\xi '\| < \|\xi \|$?  First of
all, in the right hand side of \eqref{tfm} there are no submodular
symbols of the form $[v_{i},v_{j}]$.  In fact, any submodular symbol
involving a point $v_{i}$ also includes a reducing point for 
$[v_{i},v_{j}]$.

On the other hand, consider the submodular symbols in \eqref{tfm} of
the form $[w_{i},w_{j}]$.  Since there is no relationship among the
$w_{i}$, one has no reason to believe that these modular symbols are
closer to unimodularity than those in $\uu $.  Indeed, for certain
choices of reducing points it can happen that $\|[w_{i},w_{j}]\|\geq
\|\uu \|$.

The upshot is that some care must be taken in choosing reducing
points.  In \cite[Conjectures 3.5 and 3.6]{experimental} we describe
two methods for finding reducing points for modular symbols, one using
\Vor reduction and one using LLL-reduction.  Our experience is that
if one selects reducing points using either of these conjectures,
then $\|[w_{i},w_{j}]\|< \|\uu \|$ for each of the new modular symbols
$[w_{i},w_{j}]$.  In fact, in practice these symbols are trivial or
satisfy $\|[w_{i},w_{j}]\|=1$.

\subsection{}\label{sl2.stop}
%
In the previous discussion we assumed that no submodular symbols of
any $\uu \in \support \xi $ were unimodular.  Now we say what to
do if some are.  There are three cases to consider.  

First, all submodular symbols of $\uu $ may be unimodular.  In this
case there are no reducing points, and \eqref{tfm} becomes
\begin{equation}\label{tfm2}
[v_{1},v_{2},v_{3}]\longmapsto [v_{1},v_{2},v_{3}].
\end{equation}

Second, one submodular symbol of $\uu $ may be nonunimodular, say the
symbol $[v_{1},v_{2}]$.  In this case, to build $\eta$, we use a
tetrahedron $P'$
and put $\eta (\uu ) = [v_{1},v_{2},v_{3},w_{3}]$
(Figure~\ref{tetra.fig}).  Since $[v_{1}, v_{2}, w_{3}]$ vanishes
in the boundary of $\eta $ mod $\Gamma $, \eqref{tfm} becomes
\begin{equation}\label{tfm3}
[v_{1},v_{2},v_{3}]\mapsto -[v_{1},v_{3},w_{3}] + [v_{2}, v_{3},
w_{3}].
\end{equation}

\begin{figure}[ht]
\begin{center}
\psfrag{v1}{$v_{1}$}
\psfrag{v2}{$v_{2}$}
\psfrag{v3}{$v_{3}$}
\psfrag{w3}{$w_{3}$}
\includegraphics[scale=.5]{diagrams/tetra}
\end{center}
\caption{\label{tetra.fig}}
\end{figure}

Finally, two submodular symbols of $\uu $ may be nonunimodular, say
$[v_{1},v_{2}]$ and $[v_{1},v_{3}]$.  In this case we use the cone on
a square $P''$ (Figure~\ref{cone-on-square.fig}).  To construct $\eta
(\uu )$, we must choose a decomposition of $P''$ into tetrahedra.
Since $P''$ has a nonsimplicial face, this choice affects $\xi '$ (in
contrast to the previous cases).  If we subdivide $P''$ by connecting
the vertex labelled $v_{2}$ with the vertex labelled $w_{2}$, we
obtain
\begin{equation}\label{tfm4}
[v_{1},v_{2},v_{3}]\longmapsto[v_{2},w_{2},w_{3}]+[v_{2},v_{3},w_{2}]+
[v_{1},v_{3},w_{2}].
\end{equation}

\begin{figure}[ht]
\begin{center}
\psfrag{v1}{$v_{1}$}
\psfrag{v2}{$v_{2}$}
\psfrag{v3}{$v_{3}$}
\psfrag{w2}{$w_{2}$}
\psfrag{w3}{$w_{3}$}
\includegraphics[scale=.5]{diagrams/cone-on-square}
\end{center}
\caption{\label{cone-on-square.fig}}
\end{figure}

\subsection{}
%
Now consider general $n$.  The basic technique is the same, but the
combinatorics become more complicated.  Suppose $\uu = [v_{1},\dots ,v_{n+1}]$
satisfies $q (\uu )\not =0$ in a $1$-sharbly cycle $\xi$, and for
$i=1,\dots ,n+1$ let $\vv_{i} $ be the submodular symbol $[v_{1},\dots
,\widehat{v_{i}},\dots ,v_{n+1}]$.  Assume that all $\vv_{i}$ are
nonunimodular, and for each $i$ let $w_{i}$ be a reducing point for
$\vv_{i}$.

For any subset $I\subset\{1,\dots ,n+1 \}$, let $\uu_{I}$ be the
$1$-sharbly $[u_{1},\dots ,u_{n+1}]$, where $u_{i}=w_{i}$ if $i\in I$,
and $u_{i}=v_{i}$ otherwise.  The polytope $P$ used to build $\eta
(\uu)$ is the \defn{cross polytope}, which is the higher-dimensional
analogue of the octahedron \cite[\S4.4]{experimental}.  We suppress the
details and give the final answer:
\eqref{tfm} becomes
\begin{equation}\label{tfmrelation}
\uu \longmapsto -\sum _{I} (-1)^{\#I}\uu_{I},
\end{equation}
where the sum is taken over all subsets $I\subset \{1,\dotsc ,n+1 \}$ 
of cardinality \emph{at least $2$}.  

More generally, if some $\vv_{i}$ happen to be unimodular, then the
polytope used to build $\eta$ is an iterated cone on a
lower-dimensional cross polytope.  This is already visible for $n=2$:
\begin{itemize}
\item The 2-dimensional cross polytope is a
square, and the polytope $P''$ is
a cone on a square.
\item The 1-dimensional cross polytope is an interval, and the
polytope $P'$ is a double cone on an interval.
\end{itemize}
Altogether there are $n+1$ relations generalizing
\eqref{tfm2}--\eqref{tfm4}.

\subsection{}\label{ss:sl4computations}
%
Now we describe how these computations are carried out in practice,
focusing on $\Gamma = \Gamma_{0} (N)\subset \SL_{4} (\Z)$ and $H^{5}
(\Gamma; \C)$.  Besides discussing technical details,
we also have to slightly modify some aspects of the construction in
Section~\ref{sl2.start}, since $\Gamma$ is not torsion-free.

Let $W$ be the well-rounded retract.  We can represent
a cohomology class $\beta \in H^{5} (\Gamma ; \C)$ as $ \beta = \sum q
(\sigma )\sigma $, where $\sigma $ denotes a codimension $1$ cell in
$W$.  In this case there are three types of codimension $1$ cells in
$W$.  Under the bijection $W \leftrightarrow \sC$, these cells
correspond to the \Vor cells indexed by the columns of the matrices 
\begin{equation}\label{eq:4onesharbs}
\left(\begin{array}{ccccc}
1&0&0&0&1\\
0&1&0&0&1\\
0&0&1&0&1\\
0&0&0&1&1
\end{array} \right),\quad 
\left(\begin{array}{ccccc}
1&0&0&0&1\\
0&1&0&0&1\\
0&0&1&0&1\\
0&0&0&1&0
\end{array} \right),\quad 
\left(\begin{array}{ccccc}
1&0&0&0&1\\
0&1&0&0&1\\
0&0&1&0&0\\
0&0&0&1&0
\end{array} \right).
\end{equation}
Thus each $\sigma$ in $W$ modulo $\Gamma$ corresponds to an $\SL_{4}
(\Z)$-translate of one of the matrices in \eqref{eq:4onesharbs}.  
These translates
determine basis $1$-sharblies $\uu$ (by taking the points $u_{i}$ to
be the columns), and hence we can represent $\beta$ by a 1-sharbly
chain $ \xi = \sum q
(\uu) \uu \in S_{1}$ that is a cycle in the complex of coinvariants
$(S_{*,\Gamma}, \partial_{\Gamma})$.

To make later computations more efficient, we precompute more data
attached to $\xi$.  Given a $1$-sharbly $\uu = [u_{1},\dotsc
,u_{n+1}]$, a \emph{lift} $M (\uu)$ of $\uu$ is defined to be an
integral matrix with primitive columns $M_{i}$ such that $\uu =
[M_{1},\dotsc ,M_{n+1}]$.  Then we encode $\xi $, once and for all, by
a finite collection $\Phi$ of $4$-tuples 
\[
(\uu , n (\uu ), \{\vv \}, \{M (\vv
) \}),
\]
where
\begin{enumerate}
\item $\uu$ ranges over the support of $\xi$,
\item $n (\uu )\in \C $ is the coefficient of $\uu$ in $\xi$,
\item $\{\vv \}$ is the set of submodular symbols appearing in the
boundary of $\uu$, and 
\item $\{M (\vv ) \}$ is a set of lifts for $\{\vv \}$.\label{ifour}  
\end{enumerate}
Moreover, the lifts in (\ref{ifour}) are chosen to satisfy the
following $\Gamma $-equivariance condition.  Suppose that for $\uu
,\uu '\in \support \xi $ we have $\vv \in \support (\partial \uu ) $
and $\vv '\in \support (\partial \uu ')$ satisfying $\vv = \gamma
\cdot \vv '$ for some $\gamma \in \Gamma $.  Then we require $M (\vv
)= \gamma M (\vv ')$.  This is possible since $\xi$ is a cycle modulo
$\Gamma$, although there is one complication since $\Gamma$ has torsion:
it can happen that some submodular symbol
$\vv$ of a $1$-sharbly $\uu$ is identified to \emph{itself} by an
element of $\Gamma$.  
This means that in constructing $\{M (\vv)
\}$ for $\uu$, we must somehow choose more than one lift for $\vv$.
To deal with this, let $M
(\vv)$ be any lift of $\vv$, and let $\Gamma (\vv)\subset \Gamma $ be
the stabilizer of $\vv$.  Then in $\xi$, we replace $q (\uu ) \uu$ by
\[
\frac{1}{\#\Gamma (\vv )}\sum_{\gamma \in \Gamma (\vv)} q (\uu) \uu_{\gamma},
\]
where $\uu_{\gamma }$ has the same data as $\uu$, 
except\footnote{In fact, we can be slightly
more clever than this and only introduce denominators 
that are powers of $2$.} 
that we give
$\vv$ the lift $\gamma M (\vv)$.


Next we compute and store the 1-sharbly transformation laws generalizing
\eqref{tfm2}--\eqref{tfm4}.  As a part of this we fix triangulations
of certain cross polytopes as in \eqref{tfm4}.  

We are now ready to begin the actual reduction algorithm.  We take a
Hecke operator $T (l,k)$ and build the coset representatives $\Omega $
as in \eqref{how.hecke.acts}.  For each $h\in \Omega $ and each
$1$-sharbly $\uu$ in the support of $\xi$, we obtain a non-reduced
$1$-sharbly $\uu_{h} := h\cdot\uu$.  Here $h$ acts on all the data
attached to $\uu$ in the list $\Phi$.  In particular, we replace each
lift $M (\vv)$ with $h\cdot M (\vv)$, where the dot means matrix
multiplication.

Now we check the submodular symbols of $\uu_{h}$ and choose reducing
points for the nonunimodular symbols.  This is where the lifts come in
handy.  Recall that reduction points must be chosen
$\Gamma$-equivariantly over the entire cycle.  Instead of explicitly
keeping track of the identifications between modular symbols, we do
the following trick:
\begin{enumerate}
\item Construct the \defn{Hermite normal form} $M_{\her}(\vv)$ of the lift $M
(\vv)$ (see \cite[\S2.4]{cohen:course_ant} and \exref{ch:linalg}{ex:hnf}).  Record the transformation matrix $U \in \GL_{4} (\Z)$ such
that $UM(\vv ) = M_{\her} (\vv)$.
\item Choose a reducing point $u$ for $M_{\her} (\vv)$.
\item Then the reducing point for $M (\vv)$ is $U^{-1}u$.
\end{enumerate}
This guarantees $\Gamma$-equivariance: if $\vv$, $\vv '$ are submodular
symbols of $\xi$ with $\gamma \cdot \vv = \vv '$ and with reducing
points $u, u'$, we have $\gamma u = u'$.  The reason is that the
Hermite normal form $M_{\her} (\vv)$ is a \emph{uniquely determined}
representative of the $\GL_{4}(\Z)$-orbit of
$M (\vv)$ \cite{cohen:course_ant}.  Hence if $\gamma M
(\vv) = M (\vv ')$, then $M_{\her} (\vv) = M_{\her} (\vv
')$.

After computing all reducing points, we apply the appropriate
transformation law.  The result will be a chain of $1$-sharblies, each
of which has (conjecturally) smaller norm than the original
$1$-sharbly $\uu$.  We output these $1$-sharblies if they are reduced;
otherwise they are fed into the reduction algorithm again.  Eventually
we obtain a reduced $1$-sharbly cycle $\xi '$ homologous to the
original cycle $\xi$.

The final step of the algorithm is to rewrite $\xi '$ as a cocycle on
$W$.  This is easy to do since the relevant cells of $W$ are in
bijection with the reduced $1$-sharblies.  There are some nuisances in
keeping orientations straight, but the computation is not difficult.
We refer to \cite{computation} for details.

\subsection{}
%
We now give some examples, taken from \cite{computation}, of Hecke
eigenclasses in $H^{5} (\Gamma_{0} (N); \C)$ for various levels
$N$.
Instead of giving a table of eigenvalues, we give the \defn{Hecke
polynomials}.  If $\beta$ is an eigenclass with $T (l,k) (\beta) = a
(l,k)\beta$, then we define
\[
H (\beta , l) = \sum_k (-1)^k l^{k(k-1)/2} a(l, k) X^k \in \C [X].
\]
For almost all $l$, after putting $X = l^{-s}$ where $s$ is a complex
variable, the function $H
(\beta , s)$ is the inverse of the local factor at $l$ of the
automorphic representation attached to $\beta$.
\begin{example}
Suppose $N=11$.  Then the cohomology $H^{5} (\Gamma_{0} (11); \C)$ is
2-dimensional.  There are two Hecke eigenclasses $u_{1}, u_{2}$, each
with rational Hecke eigenvalues.

\medskip
\begin{center}
\begin{tabular}{|p{40pt}|p{40pt}|p{200pt}|}
\hline
$u_{1}$&$T_2$&$ (1-4 X)(1-8 X)(1 +2 X + 2 X^2)$\\
&$T_3$&$ (1-9 X)(1-27 X)(1 + X + 3 X^2) $\\
&$T_5$&$ (1-25 X)(1-125 X)(1 - X + 5 X^2) $\\
&$T_7$&$ (1-49 X)(1-343 X)(1 +2 X + 7 X^2) $\\
\hline
$u_{2}$&$T_2$&$ (1- X)(1-2 X)(1 +8 X + 32 X^2)$\\
&$T_3$&$ (1- X)(1-3 X)(1 +9 X + 243 X^2) $\\
&$T_5$&$ (1- X)(1-5 X)(1 -25 X + 3125 X^2) $\\
&$T_7$&$ (1- X)(1-7 X)(1 +98 X + 16807 X^2) $\\
\hline
\end{tabular}
\end{center}

  
\end{example}
 
\begin{example}
Suppose $N=19$.  Then the cohomology $H^{5}
(\Gamma_{0} (19); \C)$ is 3-dimensional.  There are three Hecke
eigenclasses $u_{1}, u_{2}, u_{3}$, each with rational Hecke eigenvalues.

\medskip
\begin{center}
\begin{tabular}{|p{40pt}|p{40pt}|p{200pt}|}
\hline
$u_{1}$&$T_2$&$ (1-4 X)(1-8 X)(1  + 2 X^2)$\\
&$T_3$&$ (1-9 X)(1-27 X)(1 +2 X + 3 X^2) $\\
&$T_5$&$ (1-25 X)(1-125 X)(1 -3 X + 5 X^2) $\\
\hline
$u_{2}$&$T_2$&$ (1- X)(1-2 X)(1  + 32 X^2)$\\
&$T_3$&$ (1- X)(1-3 X)(1 +18 X + 243 X^2) $\\
&$T_5$&$ (1- X)(1-5 X)(1 -75 X + 3125 X^2) $\\
\hline
$u_{3}$&$T_2$&$ (1- 2 X)(1-4 X)(1 +3 X + 8 X^2)$\\
&$T_3$&$ (1- 3 X)(1-9 X)(1 +5 X + 27 X^2) $\\
&$T_5$&$ (1- 5 X)(1-25 X)(1 +12 X + 125 X^2) $\\
\hline
\end{tabular}
\end{center}

\end{example}

In these examples, the cohomology is completely accounted for by the
Eisenstein summand of \eqref{franke.decomp}.  In fact, let
$\Gamma'_{0} (N)\subset \SL_{2} (\Z)$ be the usual Hecke congruence
subgroup of matrices upper-triangular modulo $N$.  Then the cohomology
classes above actually come from classes in $H^{1} (\Gamma_{0}' (N))$,
that is from holomorphic modular forms of level $N$.

For $N=11$, the space of weight two cusp forms $S_{2} (11)$ is
1-dimensional.  This cusp form $f$ lifts in two different ways to
$H^{5} (\Gamma_{0} (11); \C)$, which can be seen from the quadratic
part of the Hecke polynomials for the $u_i$.  Indeed, for $u_{i}$ the
quadratic part is exactly the inverse of the local factor for the
$L$-function attached to $f$, after the substitution $X = l^{-s}$.
For $u_{2}$, we see that the lift is also twisted by the square of the
cyclotomic character.  (In fact the linear terms of the Hecke
polynomials come from powers of the cyclotomic character.)

For $N=19$, the space of weight two cusp forms $S_{2} (19)$ is again
1-dimensional.  The classes $u_{1}$ and $u_{2}$ are lifts of this
form, exactly as for $N=11$.  The class $u_{3}$, on the other hand,
comes from $S_{4} (19)$, the space of weight $4$ cusp forms on
$\Gamma_{0}' (19)$.  In fact, $\dim S_{4} (19) = 4$, with one Hecke
eigenform defined over $\Q$ and another defined over a totally real
cubic extension of $\Q$.  Only the rational weight four eigenform
contributes to $H^{5} (\Gamma_{0} (19); \C)$.  One can show that
whether or not a weight four cuspidal eigenform $f$ contributes to the
cohomology of $\Gamma_{0} (N)$ depends only on the sign of the
functional equation of $L (f,s)$ \cite{weselman}.  This phenomenon
is typical of what one encounters when studying Eisenstein cohomology.

In addition to the lifts of weight 2 and weight 4 cusp forms, for
other levels one finds lifts of Eisenstein series of weights 2 and 4
and lifts of cuspidal cohomology classes from subgroups of $\SL_{3}
(\Z)$.  For some levels one finds cuspidal classes that appear to be
lifts from the group of symplectic similitudes $\GSp (4)$.  More
details can be found in \cite{computation, computationII}.

\subsection{}
%
Here are some notes on the reduction algorithm and its implementation:
\begin{itemize}
\item Some additional care must be taken when selecting reducing
points for the submodular symbols of $\uu$.  In particular, in
practice one should choose $w$ for $\vv$ such that $\sum \| \vv_{i}
(w)\|$ is minimized.  Similar remarks apply when choosing a
subdivision of the crosspolytopes in Section~\ref{sl2.stop}.
\item In practice, the reduction algorithm has \emph{always} terminated with
a reduced $1$-sharbly cycle $\xi '$ homologous to $\xi$.  However, at
the moment we cannot prove that this will always happen.
\item Experimentally, the efficiency of the reduction step appears to
be comparable to that of Theorem \ref{thm:modsymbalg}.  In other words
the depth of the ``reduction tree'' associated to a given $1$-sharbly
$\uu$ seems to be bounded by a polynomial in $\log \log \| \uu \| $.  Hence
computing the Hecke action using this algorithm is extremely
efficient.

On the other hand, computing Hecke operators on $\SL_{4}$ is still a
much bigger computation---relative to the level---than on $\SL_{2}$
and $\SL_{3}$.  For example, the size of the full retract $W$ modulo
$\Gamma_{0} (p)$ is roughly $O (p^{6})$, which grows rapidly with $p$.
The portion of the retract corresponding to $H^{5}$ is much smaller,
around $p^{3}/10$, but this still grows quite quickly.  This makes
computing with $p>100$ out of reach at the moment.

The number of Hecke cosets grows rapidly as well, e.g., the number of
coset representatives of $T (l,2)$ is $l^{4}+l^{3}+2l^{2}+l+1$.  Hence
it is only feasible to compute Hecke operators for small $l$; for
large levels only $l=2$ is possible.

Here are some numbers to give an idea of the size of these
computations.  For level $73$, the rank of $H^{5}$ is 20.  There are
39504 cells of codimension~$1$ and 4128 top-dimensional cells in $W$
modulo $\Gamma_{0} (73)$.  The computational techniques in
\cite{computation} used at this level (a Lanczos scheme over a large
finite field) tend to produce sharbly cycles supported on almost all
the cells.  Computing $T (2,1)$ requires a reduction tree of 
depth $1$
and produces as many as 26 reduced $1$-sharblies for each of the 15
nonreduced Hecke images.  Thus one cycle produces a cycle supported on
as many as 15406560 $1$-sharblies, all of which must be converted to
an appropriate cell of $W$ modulo $\Gamma$.  Also this is just what
needs to be done for \emph{one} cycle; do not forget that the rank of
$H^{5}$ is 20.

In practice the numbers are slightly better, since the reduction step
produces fewer $1$-sharblies on average and since the support of the
initial cycle has size less than $39504$.  Nevertheless the orders of
magnitude are correct.  

\item Using lifts is a convenient way to encode the global
$\Gamma$-identifications in the cycle $\xi$, since it means we do not
have to maintain a big data structure keeping track of the identifications on
$\partial \xi$.  However, there is a certain expense in computing the
Hermite normal form.  This is balanced by the benefit that working
with the data $\Phi$ associated to $\xi$ allows us to reduce the
supporting $1$-sharblies $\uu$ \emph{independently}.  This means we
can cheaply parallelize our computation: each $1$-sharbly, encoded as
a $4$-tuple $(\uu , n (\uu ), \{\vv \}, \{M (\vv ) \})$, can be
handled by a separate computer.  The results of all these individual
computations can then be collated at the end, when producing a
$W$-cocycle.
\end{itemize}

%%%%%%%%%%%%%%%%%%%
%                 %
%  open problems  %
%                 %
%%%%%%%%%%%%%%%%%%%

\section{Complements and Open Problems}\label{apps}

\subsection{}
%
We conclude this appendix by giving some complements and describing
some possible directions for future work, both theoretical and
computational.  Since a full explanation of the material in this
section would involve many more pages, we will be brief and will
provide many references.

\subsection{Perfect Quadratic Forms over Number Fields and Retracts}
%
Since \Vor 's pioneering work \cite{vor1}, it has been the goal of
many to extend his results from $\Q$ to a general algebraic number
field $F$.  Recently Coulangeon \cite{coul}, building on work of Icaza
and Baeza \cite{icaza, baeza.icaza}, has found a good notion of
\defn{perfection} for quadratic forms over number
fields\footnote{Such forms are called \defn{Humbert forms} in the
literature.}.  One of the key ideas in \cite{coul} is that the correct
notion of equivalence between Humbert forms involves not only the
action of $\GL_{n} (\sO_{F})$, where $\sO_{F}$ is the ring of integers
of $F$, but also the action of a certain continuous group $U$ related
to the units $\sO_{F}^{\times }$.  One of Coulangeon's basic results
is that there are finitely many equivalence classes of perfect Humbert
forms modulo these actions.

On the other hand, Ash's original construction of retracts
\cite{ash77} introduces a geometric notion of perfection.  Namely he
generalizes the \Vor polyhedron $\Pi$ and defines a quadratic form to
be perfect if it naturally indexes a facet of $\Pi$.  What is the
connection between these two notions?  Can one use Coulangeon's
results to construct cell complexes to be used in cohomology
computations?  One tempting possibility is to try to use the group $U$ 
to collapse the \Vor cells of \cite{ash77} into a
cell decomposition of the symmetric space associated to $\SL_{n} (F)$.

\subsection{The Modular Complex}
%
In his study of multiple $\zeta$-values, Goncharov has recently defined
the \defn{modular complex} $M^{*}$ \cite{gonch1, gonch2}.  This is an
$n$-step complex of $\GL_{n} (\Z)$-modules closely related both to the
properties of multiple polylogarithms evaluated at $\mu_{N}$, the
$N$th roots of unity, and to the action of $G_{\Q}$ on $\pi_{1,N} =
\pi_{1}^{l} (\Proj^{1}\smallsetminus \{0,\infty ,\mu_{N} \})$, the
pro-$l$ completion of the algebraic fundamental group of
$\Proj^{1}\smallsetminus \{0,\infty ,\mu_{N} \}$.  

Remarkably, the modular complex is very closely related to the \Vor
decomposition $\sV$.  In fact, one can succinctly describe the modular
complex by saying that it is the chain complex of the cells coming
from the top-dimensional \Vor cone of type $A_{n}$.  This is all of
the \Vor decomposition for $n=2,3$, and Goncharov showed that the
modular complex is quasi-isomorphic to the full \Vor complex for
$n=4$.  Hence there is a precise relationship among multiple
polylogarithms, the Galois action on $\pi_{1,N}$, and the cohomology
of level $N$ congruence subgroups of $\SL_{n} (\Z)$.

The question then arises, how much of the cohomology of congruence
subgroups is captured by the modular complex for all $n$?  Table
\ref{vortable} indicates that asymptotically very little of the \Vor
decomposition comes from the $A_n$ cone, but this says nothing about
the cohomology.  The first interesting case to consider is $n=5$.


\subsection{Retracts for Other Groups}
%
The most general construction of retracts $W$ known \cite{ash.wrr}
applies only to \emph{linear} symmetric spaces.\index{linear symmetric spaces}  
The most familiar
example of such a space is $\SL_{n} (\R)/\SO (n)$; other examples are
the symmetric spaces associated to $\SL_{n}$ over number fields and
division algebras.

Now let $\Gamma\subset \bG (\Q )$ be an arithmetic group, and let $X =
G/K$ be the associated symmetric space.  What can one say about cell
complexes that can be used to compute $H^{*} (\Gamma; \MMM)$?  The
theorem of Borel--Serre mentioned in Section~\ref{ss:vcd} implies the
vanishing of $H^{k}(\Gamma; \MMM )$ for $k>\nu := \dim X - q$, where
$q$ is the \defn{\protect{$\Q$-rank}} of $\Gamma$.  For example, for the split
form of $\SL_{n}$, the $\Q$-rank is $n-1$.  For the split symplectic
group $\SP_{2n}$, the $\Q$-rank is $n$.  Moreover, this bound is
sharp: there will be coefficient modules $\MMM$ for which $H^{\nu}
(\Gamma ; \MMM) \not = 0$.  Hence any minimal cell complex used to compute the
cohomology of $\Gamma$ should have dimension $\nu$. 

Ideally one would like to see such a complex realized as a subspace of~$X$ 
and would like to be able to treat all finite index subgroups of
$\Gamma $ simultaneously.  This leads to the following question: is
there a $\Gamma$-equivariant deformation retraction of $X$ onto a
regular cell complex $W$ of dimension $\nu$?

For $\bG = \SP_{4}$, McConnell and MacPherson showed that the answer is
yes.  Their construction begins by realizing the symplectic symmetric
space $X_{\SP}$ as a subspace of the special linear symmetric space
$X_{\SL}$.  They then construct subsets of $X_{\SP}$ by intersecting
the \Vor cells in $X_{\SL}$ with $X_{\SP}$.  Through explicit
computations in coordinates they prove that these intersections are
cells and give a cell decomposition of $X_{\SP}$.  By taking an
appropriate dual complex (as suggested by Figures \ref{tess} and
\ref{retract2.fig} and as done in \cite{ash77}), they construct the
desired cell complex $W$.

Other progress has been recently made by Bullock \cite{bullock},
Bullock and Connell \cite{bullock.connell}, and Yasaki \cite{yasaki1,
yasaki2} in the case of groups of $\Q$-rank 1.  In particular, Yasaki
uses the \defn{tilings} of Saper \cite{saper} to construct an explicit
retract for the unitary group $\SU (2,1)$ over the Gaussian integers.
His method also works for Hilbert modular groups, although further
refinement may be needed to produce a regular cell complex.  Can one
generalize these techniques to construct retracts for groups of
arbitrary $\Q$-rank?  Is there an analogue of the \Vor decomposition
for these retracts (i.e., a dual cell decomposition of the symmetric
space)?  If so, can one generalize ideas in
Sections~\ref{mod.symb.section}--\ref{other.coho} and use 
that generalization to compute the
action of the Hecke operators on the cohomology?



\subsection{Deeper Cohomology Groups}
%
The algorithm in Section~\ref{other.coho} can be used
to compute the Hecke action on
$H^{\nu -1} (\Gamma)$.  For $n>4$, this group no longer contains
cuspidal cohomology classes.  Can one generalize this algorithm to
compute the Hecke action on deeper cohomology groups?  The first
practical case is $n=5$.  Here $\nu = 10$, and the highest degree in
which cuspidal cohomology can live is $8$.  This case is also
interesting since the cohomology of full level has been studied
\cite{gangl}.

Here are some indications of what one can expect.  The general
strategy is the same: for a $k$-sharbly $\xi$ representing a
class in $H^{\nu -k} (\Gamma)$, begin by $\Gamma$-equivariantly
choosing reducing points for the nonunimodular submodular symbols of
$\xi$.  This data can be packaged into a new $k$-sharbly cycle as in
Section~\ref{ss:buildingxiprime}ff, but the crosspolytopes must be replaced
with \defn{hypersimplices}.  By definition, the hypersimplex $\Delta
(n,k)$ is the convex hull in $\R^{n}$ of the points $\{\sum_{i\in I}
e_{i} \}$, where $I$ ranges over all order $k$ subsets of $\{1,\dotsc
,n \}$ and $e_{1},\dotsc ,e_{n}$ denotes the standard 
basis of $\R ^{n}$.

The simplest example is $n=2$, $k=2$.  From the point of view of
cohomology, this is even less interesting than $n=2$, $k=1$, since now
we are computing the Hecke action on $H^{-1} (\Gamma)$!  Nevertheless,
the geometry here illustrates what one can expect in general.

Each $2$-sharbly in the support of $\xi$ can be written as $
[v_{1},v_{2},v_{3},v_{4}]$ and determines six submodular symbols, of
the form $[v_{i}, v_{j}]$, $i\not =j$.  Assume for simplicity that all these
submodular symbols are nonunimodular.  Let $w_{ij}$ be the reducing
point for $[v_{i},v_{j}]$.  Then use the ten points $v_{i}, w_{ij}$ to
label the vertices of the hypersimplex $\Delta (5,2)$ as in Figure
\ref{hy1simp.fig} (note that $\Delta (5,2)$ is $4$-dimensional).

\begin{figure}[htb]
\psfrag{04}{$v_{1}$}
\psfrag{14}{$v_{2}$}
\psfrag{24}{$v_{3}$}
\psfrag{34}{$v_{4}$}
\psfrag{01}{$w_{12}$}
\psfrag{02}{$w_{13}$}
\psfrag{03}{$w_{14}$}
\psfrag{12}{$w_{23}$}
\psfrag{13}{$w_{24}$}
\psfrag{23}{$w_{34}$}
\begin{center}
\includegraphics[scale=0.4]{diagrams/hy1simp_2}
\end{center}
\caption{\label{hy1simp.fig}}
\end{figure}


The boundary of this hypersimplex gives the analogue of
\eqref{tfm}.  Which $2$-sharblies will appear in $\xi '$?  The boundary
$\partial \Delta (5,2)$ is a union of five tetrahedra and five octahedra.
The outer tetrahedron will not appear in $\xi '$, since that is the
analogue of the left side of \eqref{tfm}.  The four octahedra sharing a
triangular face with the outer tetrahedron also will not appear, since
they disappear when considering $\xi '$ modulo $\Gamma$.  The
remaining four tetrahedra and the central octahedron survive to $\xi '$
and constitute the right side of the analogue of \eqref{tfm}.  Note that we
must choose a simplicial subdivision of the central octahedron to write
the result as a $2$-sharbly cycle and that this must be done with
care since it introduces a new submodular symbol.

If some submodular symbols are unimodular, then again one must
consider iterated cones on hypersimplices, just as in
Section~\ref{sl2.stop}.  The analogues of these steps become more
complicated, since there are now many simplicial subdivisions of a
hypersimplex\footnote{Indeed, computing all simplicial subdivisions of
$\Delta (n,k)$ is a difficult problem in convex geometry.}. There is
one final complication: in general we cannot use reduced $k$-sharblies
alone to represent cohomology classes.  Thus one must terminate the
algorithm when $\|\xi \|$ is less than some predetermined bound.

\subsection{Other Linear Groups}
%
Let $F$ be a number field, and let $\bG = \bR_{F/\Q} (\SL_{n})$
(Example \ref{ex:resofscalars}).  Let $\Gamma \subset \bG (\Q)$ be an
arithmetic subgroup.  Can one compute the action of the
Hecke operators on $H^{*} (\Gamma)$?  

There are two completely different approaches to this problem.  The
first involves the generalization of the modular symbols method.  One
can define the analogue of the sharbly complex, and can try to extend
the techniques of Sections~\ref{mod.symb.section}--\ref{other.coho}.

This technique has been extensively used when 
$F$ is \emph{imaginary quadratic} and $n=2$.  We have $X = \SL_{2}
(\C)/\SU (2)$, which is isomorphic to $3$-dimensional hyperbolic
space $\fh_{3}$.  The arithmetic groups $\Gamma \subset \SL_{2}
(\sO_{F})$ are known as \defn{Bianchi groups}.  The retracts and
cohomology of these groups have been well studied; as a representative
sample of works we mention \cite{mendoza, egm, vogtmann,
grune.sch}.

Such groups have $\Q$-rank 1 and thus have cohomological dimension
$2$.  One can show that the cuspidal classes live in degrees $1$ and
$2$.  This means that we can use modular symbols to investigate the
Hecke action on cuspidal cohomology.  This was done by Cremona
\cite{crem} for \emph{euclidean} fields $F$.  In that case Theorem
\ref{thm:modsymbalg} works with no trouble (the euclidean algorithm is
needed to construct reducing points).  For noneuclidean fields
further work has been done by Whitley \cite{whit}, Cremona and Whitely
\cite{crem.whit} (both for principal ideal domains), Bygott
\cite{bygott} (for $F=\Q (\sqrt{-5})$ and any field with class group
an elementary abelian $2$-group), and Lingham \cite{lingham} (any
field with odd class number).  Putting all these ideas together allows
one to generalize the modular symbols method to \emph{any} imaginary
quadratic field \cite{crem.letter}.

For $F$ imaginary quadratic and $n>2$, very little has been studied.
The only related work to the best of our knowledge is that of
Staffeldt \cite{staffeldt}.  He determined the structure of the \Vor
polyhedron in detail for $\bR_{F/\Q} (\SL_{3})$, where $F = \Q
(\sqrt{-1})$.  We have $\dim X = 8$ and $\nu =6$.  The cuspidal
cohomology appears in degrees $3,4,5$, so one could try to use the
techniques of Section~\ref{other.coho} to investigate it.

Similar remarks apply to $F$ \emph{real quadratic} and $n=2$.  The
symmetric space $X \simeq \fh \times \fh$ has dimension $4$ and the
$\Q$-rank is 1, which means $\nu = 3$.  Unfortunately the cuspidal
cohomology appears only in degree $2$, which means modular symbols
cannot see it.  On the other hand, 1-sharblies can see it, and so one
can try to use ideas in Section~\ref{other.coho} here to compute the Hecke
operators.  The data needed to build the retract $W$ already
(essentially) appears in the literature for certain fields; see for
example \cite{ong}.

The second approach shifts the emphasis from modular symbols and the
sharbly complex to the \Vor fan and its cones.  For this approach we
must assume that the group $\Gamma$ is associated to a
\defn{self-adjoint homogeneous cone} over $\Q$.  (cf. \cite{ash77}).
This class of groups includes arithmetic subgroups of $\bR_{F/\Q}
(\SL_{n})$, where $F$ is a totally real or CM field.  Such groups have
all the nice structures in Section~\ref{vcd.section}.  For example, we have
a cone $C$ with a $G$-action.  We also have an analogue of the \Vor
polyhedron $\Pi$.  There is a natural compactification $\tilde{C}$ of
$C$ obtained by adjoining certain self-adjoint homogeneous cones of
lower rank.  The quotient $\Gamma \backslash \tilde{C}$ is singular in
general, but it can still be used to compute $H^{*} (\Gamma ; \C)$.  The
polyhedron $\Pi$ can be used to construct a fan $\sV$ that gives a
$\Gamma$-equivariant decomposition of all of $\tilde{C}$.  But the
most important structure we have is the \Vor reduction algorithm:
given any point $x\in \tilde{C}$, we can determine the unique \Vor
cone containing $x$.

Here is how this setup can be used to compute the Hecke action.  Full
details are in \cite{msa, ccalg}.  We define two chain complexes
$\bC^{V}_{*}$ and $\bC^{R}_{*}$.  The latter is essentially the chain
complex generated by all simplicial rational polyhedral cones in
$\tilde{C}$; the former is the subcomplex generated by the \Vor cones.
These are the analogues of the sharbly complex and the chain complex
associated to the retract $W$, and one can show that either can be used
to compute $H^{*} (\Gamma ; \C)$.  Take a cycle $\xi \in \bC^{V}_{*}$
representing a cohomology class in $H^{*} (\Gamma ; \C)$ and act on it
by a Hecke operator $T$.  We have $T (\xi) \in \bC^{R}_{*}$, and we
must push $T (\xi)$ back to $\bC^{V}_{*}$.

To do this, we use the linear structure on $\tilde{C}$ to subdivide $T
(\xi)$ very finely into a chain $\xi '$.  For each $1$-cone $\tau$ in
$\support \xi '$, we choose a $1$-cone $\rho_{\tau}\in
\tilde{C}\smallsetminus C$ and assemble them using the combinatorics
of $\xi '$ into a polyhedral chain $\xi ''$ homologous to $\xi'$.
Under certain conditions involved in the construction of $\xi '$, this
chain $\xi ''$ will lie in $\bC^{V}_{*}$.  

We illustrate this process for the split group $\SL_{2}$; more details
can be found in \cite{msa}.  We work modulo homotheties, so that the
three-dimensional cone $\tilde{C}$ becomes the extended upper
half plane $\fh^{*}:= \fh \cup \Q \cup \{\infty \} $, with $\partial
\tilde{C}$ passing to the cusps $\fh ^{*}\smallsetminus \fh $.  As
usual top-dimensional \Vor cones become the triangles of the Farey
tessellation, and the cones $\rho_{\tau}$ become cusps.  Given any
$x\in \fh$, let $R (x)$ be the set of cusps of the unique triangle or
edge containing $x$ (this can be computed using the \Vor reduction
algorithm).  Extend $R$ to a function on $\fh^{*}$ by putting $R (u) =
\{u \}$ for any cusp $u$.

In $\fh$, the support of $T (\xi)$ becomes a geodesic $\mu $ between
two cusps $u$, $u'$, in other words the support of a modular symbol
$[u,u' ]$ (Figure
\ref{partition.fig}).  Subdivide $\mu$ by choosing points
$x_{0},\dotsc ,x_{n}$ such that $x_{0}=u$, $x_{n} = u'$, and $R
(x_{i})\cap R (x_{i+1})\not =\emptyset$.  (This is easily done, for
example by repeatedly barycentrically subdividing $\mu$.)
For each $i<n$ choose a
cusp $q_{i}\in R (x_{i})\cap R (x_{i+1})$, and put $q_{n}=u'$.  Then
we have a relation in $H^{1}$:
\begin{equation}\label{eq:sahc.rel}
[u,u']=[q_{0},q_{1}]+\dotsb +[q_{n-1},q_{n}].
\end{equation}
Moreover, each $[q_{i},q_{i+1}]$ is unimodular, since $q_{i}$ and
$q_{i+1}$ are both vertices of a triangle containing $x_{i+1}$.
Upon lifting \eqref{eq:sahc.rel} back to $\bC^{R}_{*}$, the cusps
$q_{i}$ become the $1$-cones $\rho_{\tau}$ and give us a relation
$T (\xi) = \xi ''\in \bC^{V}_{*}$.

\begin{figure}[htb]
\psfrag{mu}{$\mu $}
\psfrag{u1}{$u$}
\psfrag{u2}{$u'$}
\centerline{\includegraphics[scale = .5]{diagrams/partition}}
\caption{A subdivision of $\mu $; the solid dots are the $x_{i}$.
Since the $x_{i}$ lie in the same or adjacent \Vor \ cells, we can
assign cusps to them to construct a homology to a cycle in $\bC ^{V}_{*}$.\label{partition.fig}}
\end{figure}

\subsection{The Sharbly Complex for General Groups}
%
In \cite{sympms} we generalized Theorem \ref{thm:modsymbalg} (without
the complexity statement) to the symplectic group $\SP_{2n}$.  Using
this algorithm and the symplectic retract \cite{mmc1, mmc2}, one can
compute the action of the Hecke operators on the top-degree cohomology
of subgroups of $\SP_{4} (\Z)$. 

More recently, Toth has investigated modular symbols for other groups.
He showed that the unimodular symbols generate the top-degree
cohomology groups for $\Gamma$ an arithmetic subgroup of a split
classical group or a split group of type $E_{6}$ or $E_{7}$
\cite{toth}.  His technique of proof is completely different from that
of \cite{sympms}.  In particular he does not give an analogue of the
Manin trick.  Can one extract an algorithm from Toth's proof that can
be used to explicitly compute the action of the Hecke operators on
cohomology?

The proof of the main result of \cite{sympms} uses a description of
the relations among the modular symbols.  These relations were
motivated by the structure of the cell complex in \cite{mmc1, mmc2}.
The modular symbols and these relations are analogues of the groups
$S_{0}$ and $S_{1}$ in the sharbly complex.  Can one 
extend these combinatorial constructions to form a \defn{symplectic
sharbly complex}?  What about for general groups $\bG$?  

Already for $\SP_{4}$, resolution of this question would have immediate
arithmetic applications.  Indeed, Harder has a beautiful conjecture
about certain congruences between holomorphic modular forms and Siegel
modular forms of full level \cite{harder-arbeit}.  Examples of these
congruences were checked numerically in \cite{harder-arbeit} using techniques of
\cite{faber-van.der.geer} to compute the Hecke action.

However, to investigate higher levels, one needs a different technique.
The relevant cohomology classes live in $H^{\nu -1} (\Gamma ; \MMM)$,
so one only needs to understand the first three terms of the complex
$S_{0}\leftarrow S_{1}\leftarrow S_{2}$.  We understand $S_{0}$,
$S_{1}$ from \cite{sympms}; the key is understanding $S_{2}$, which
should encode relations among elements of $S_{1}$.  If one could do
this and then could generalize the techniques of \cite{experimental},
one would have a way to investigate Harder's conjecture.

\subsection{Generalized Modular Symbols}
%
We conclude this appendix by discussing a geometric approach to
modular symbols.  This complements the algebraic approaches presented
in this book and leads to many new interesting phenomena and
problems.

Suppose $\bH$ and $\bG$ are connected semisimple algebraic groups over
$\Q$ with an injective map $f\colon \bH \rightarrow \bG$.  Let $K_{H}$
be a maximal compact subgroup of $H = \bH (\R)$, and suppose $K\subset
G$ is a maximal compact subgroup containing $f (K_{H})$.  Let $X =
G/K$ and $Y = H/K_{H}$.

Now let $\Gamma \subset \bG (\Q)$ be a torsion-free arithmetic
subgroup.  Let $\Gamma_{H} = f^{-1} (\Gamma)$.  We get a map
$\Gamma_{H}\backslash Y \rightarrow \Gamma \backslash X$, and we
denote the image by $S (H,\Gamma)$.  Any
compactly supported cohomology class $\xi \in H_{c}^{\dim Y} (\Gamma
\backslash X; \C)$ can be pulled back via $f$ to $\Gamma_{H}\backslash
Y$ and integrated to obtain a complex number.  Hence $S (H,\Gamma)$
defines a linear form on $H_{c}^{\dim Y} (\Gamma
\backslash X; \C)$.  By Poincar\'e duality, this linear form
determines a class $[S (H,\Gamma )] \in H^{\dim X - \dim Y} (\Gamma
\backslash X; \C)$, called a \defn{generalized
modular symbol}.  Such classes have been considered by many authors,
for example \cite{ash.borel, venky.speh, harder.modsym,
ash.ginzburg.rallis}.

As an example, we can take $\bG$ to be the split form of $\SL_{2}$,
and we can take $f \colon \bH \rightarrow \bG $ to be the inclusion of
connected component of the diagonal subgroup.  Hence $H \simeq
\R_{>0}$.  In this case $K_{H}$ is trivial.  The image of $Y$ in $X$
is the ideal geodesic from $0$ to $\infty$.  One way to vary $f$ is by
taking an $\SL_{2} (\Q)$-translate of this geodesic, which gives a
geodesic between two cusps.  Hence we can obtain the support of any
modular symbol this way.  This example generalizes to $\SL_{n}$ to
yield the modular symbols in Section~\ref{mod.symb.section}.  Here $H \simeq
(\R >0)^{n-1}$.  Note that $\dim Y = n-1$, so the cohomology classes
we have constructed live in the top degree $H^{\nu} (\Gamma \backslash
X; \C)$.

Another family of examples is provided by taking $\bH$ to be a Levi
factor of a parabolic subgroup; these are the modular symbols studied
in \cite{ash.borel}.

There are many natural questions to study for such objects.  Here are
two:

\begin{itemize}
\item Under what conditions on $\bG , \bH , \Gamma$ is $[S
(H,\Gamma)]$ nonzero?  This question is connected to relations between periods
of automorphic forms and functoriality lifting.  There are a variety
of partial results known; see for example \cite{venky.speh, ash.ginzburg.rallis}.
\item We know the usual modular symbols span the top-degree cohomology
for any arithmetic group $\Gamma$.  Fix a class of generalized modular
symbols by fixing the pair $\bG , \bH $ and fixing some class of maps
$f$.  How much of the cohomology can one span for a general arithmetic
group $\Gamma \subset \bG (\Q)$?

A simple example is given by the Ash--Borel construction for $\bG =
\SL_{3}$ and $\bH$ a Levi factor of a rational parabolic subgroup
$\bP$ of type $(2,1)$.  In this case $H \simeq \SL_{2} (\R) \times
\R_{>0}$ and sits inside $G$ via
\[
g\left( 
\begin{array}{cc}
\alpha^{-1}M&0\\
0&\alpha
\end{array}\right)g^{-1}, \quad M\in \SL_{2} (\R), \quad \alpha \in
\R_{>0}, \quad g \in \SL_{3} (\Q).
\]
For $\Gamma \subset \SL_{3} (\Z)$ these symbols define a subspace 
\[
S_{(2,1)}\subset H^{2} (\Gamma
\backslash X; \C).
\]
Are there $\Gamma $ for which $S_{(2,1)}$ equals
the full cohomology space?  For general $\Gamma$ how much is
captured?  Is there
a nice combinatorial way to write down the relations among these
classes?  Can one cook up a generalization of Theorem
\ref{thm:modsymbalg} for these classes and use it to compute Hecke
eigenvalues?

\end{itemize}


